\documentclass[11pt]{article}
\usepackage[T1]{fontenc}
\usepackage{lmodern}
\usepackage[margin=1in]{geometry}
\usepackage{amsmath,amssymb,amsthm,mathtools}
\usepackage{microtype}
\usepackage{enumitem,booktabs,xcolor}
\usepackage{hyperref}
\hypersetup{colorlinks=true,linkcolor=blue!45!black,citecolor=green!35!black,
  urlcolor=blue!50!black,
  pdftitle={Canonical ampleness of compact hyperbolic Kahler manifolds},
  pdfauthor={OpenAI}}
\pdftrailerid{}
\pdfsuppressptexinfo=15
\newcommand{\C}{\mathbb C}
\newcommand{\R}{\mathbb R}
\newcommand{\Z}{\mathbb Z}
\newcommand{\D}{\mathbb D}
\newcommand{\PP}{\mathbb P}
\newcommand{\dd}{\mathrm d}
\newcommand{\ddc}{i\partial\bar\partial}
\newcommand{\eps}{\varepsilon}
\newcommand{\cL}{\mathcal L}
\newcommand{\cO}{\mathcal O}
\DeclareMathOperator{\Ric}{Ric}
\DeclareMathOperator{\tr}{tr}
\DeclareMathOperator{\osc}{osc}
\newtheorem{theorem}{Theorem}[section]
\newtheorem{lemma}[theorem]{Lemma}
\newtheorem{proposition}[theorem]{Proposition}
\newtheorem{corollary}[theorem]{Corollary}
\theoremstyle{definition}
\newtheorem{definition}[theorem]{Definition}
\theoremstyle{remark}
\newtheorem{remark}[theorem]{Remark}
\numberwithin{equation}{section}
\setlist{topsep=4pt,itemsep=3pt,parsep=0pt}
\setlength{\emergencystretch}{2em}
\title{Canonical ampleness of compact hyperbolic K\"ahler manifolds}
\author{OpenAI}
\date{September 23, 2026}

\begin{document}
\maketitle
\begin{abstract}
We prove that every compact connected K\"ahler manifold of positive
complex dimension containing no nonconstant entire curve has ample
canonical bundle and is projective. This resolves positively Kobayashi's
canonical-ampleness conjecture in the smooth compact K\"ahler category.
\end{abstract}

\section{Introduction}\label{sec:introduction}

A complex manifold is \emph{Brody hyperbolic} if every holomorphic map
$\C\to X$ is constant. On a compact complex manifold, Brody's theorem
identifies this condition with Kobayashi hyperbolicity, the
nondegeneracy of the intrinsic pseudodistance defined by holomorphic
disc chains \cite{Brody1978,Kobayashi1970}. We use ``hyperbolic'' in
this sense throughout.

For a complex manifold of dimension $n$, the canonical bundle
$K_X=\bigwedge^n(T^{1,0}X)^*$ is the line bundle of holomorphic
$n$-forms. A line bundle is \emph{ample} if some positive tensor power
has global sections defining a holomorphic embedding into projective
space. Kobayashi's canonical-ampleness conjecture predicts that the
absence of nonconstant entire curves forces $K_X$ to be ample on a
compact K\"ahler manifold. It therefore connects a restriction on
one-dimensional holomorphic maps with both positivity of top-degree
forms and projective algebraicity. We prove this implication in every
positive dimension.

\begin{theorem}\label{thm:main}
Let $X$ be a compact connected K\"ahler manifold of positive complex
dimension. If every holomorphic map $\C\to X$ is constant, then $K_X$ is
ample. In particular, a positive tensor power of $K_X$ defines a
holomorphic embedding of $X$ into a complex projective space.
\end{theorem}

\subsection{Context and the main difficulty}

Kobayashi's theory supplies the intrinsic metric framework
\cite{Kobayashi1970}; the canonical-ampleness question and its
relationship with curvature are discussed in \cite{Diverio2020}.
Theorem~\ref{thm:main} gives a positive resolution in the smooth
compact K\"ahler category. A major preceding advance is the theorem
of Wu and Yau: negative holomorphic sectional curvature of a K\"ahler metric forces
canonical ampleness on a smooth projective manifold
\cite{WuYau2016}. Tosatti and Yang extended that implication to
compact K\"ahler manifolds without assuming projectivity
\cite{TosattiYang2017}. Diverio--Trapani extended the conclusion to quasi-negative holomorphic
sectional curvature: the curvature is nonpositive everywhere and
strictly negative in every nonzero tangent direction at one point
\cite{DiverioTrapani2019}. Wu--Yau subsequently gave another proof
covering both the negative and quasi-negative cases
\cite{WuYauRemark2016}. Broder--Stanfield allow the negatively curved
metric $\omega$ to be Hermitian and pluriclosed, meaning
$i\partial\bar\partial\omega=0$, while retaining the ambient
K\"ahler hypothesis \cite[Theorem~1.1]{BroderStanfield2026}.
These metric hypotheses give curvature
inequalities unavailable from the absence of entire curves alone.
The present proof begins instead with the derivative bounds supplied
by compact hyperbolicity.

A related line of work uses Gromov's K\"ahler hyperbolicity
\cite{Gromov1991}: a K\"ahler form pulls back to the exterior
derivative of a bounded one-form on the universal covering.
Building on Gromov's general-type result, Chen and Yang proved
canonical ampleness under this hypothesis
\cite[Definitions~2.2 and~2.7, Theorem~2.11]{ChenYang2018}.
This condition is distinct from Brody hyperbolicity
\cite[Corollary~4.4]{ChenYang2018}.

Recent work also establishes a weaker positivity conclusion under the
absence of rational curves. A line bundle is \emph{nef} if, for every
$\eps>0$, it admits a smooth Hermitian metric whose curvature is
bounded below by $-\eps\omega$, where $\omega$ is any fixed background
K\"ahler form. Ou's characterization of uniruled compact K\"ahler manifolds
\cite[Theorem~1.1]{Ou2025} shows that, in the absence of rational
curves, $K_X$ is pseudoeffective: its first Chern class contains a
closed positive current. Ou's theorem also supplies the
lower-dimensional hypothesis in Cao--H\"oring's rational-curve theorem
\cite[Theorem~1.3]{CaoHoring2020}. The latter then shows that failure
of nefness would produce a rational curve. Thus nefness is already
known under this weaker hypothesis. The further issue for ampleness
is positive canonical volume. We retain a direct disc proof of
nefness because it produces bounded canonical potentials and uses
the same Hessian identity as our volume estimate.

Extremal holomorphic discs have long connected intrinsic complex
geometry with analytic estimates, notably in Lempert's theory on
convex domains \cite{Lempert1981}. The extremal problem used here
includes an area penalty on discs in a compact target. Its difficulty
is at the boundary: compact convergence does not by itself justify
variations of an integral over the whole unit disc. We first obtain
boundary Sobolev regularity from maximality. We then construct a
holomorphic frame normalized on the boundary and integrate the needed
infinitesimal sections into actual holomorphic families that keep the
center fixed. These steps make one Hessian calculation available for
two different positivity estimates.

\subsection{The analytic argument}

Put $n=\dim_{\C}X$, let $\D=\{|z|<1\}$, and fix a K\"ahler metric $k$.
A disc $f:\D\to X$ has finite area if the integral of
$k(f'(z),f'(z))$ over $\D$ is finite. Compact hyperbolicity bounds the squared
$k$-norm of the central derivative of every holomorphic disc by a
constant $D$. For a closed
real $(1,1)$-form $q$, evaluated through its associated Hermitian
tensor, use the area and Poisson integrals
\[
 A_q(f)=\frac1\pi\int_\D q(f',f')\,\dd S,
 \qquad
 P_q(f)=\frac2\pi\int_\D\log\frac1{|z|}\,q(f',f')\,\dd S,
\]
where $\dd S$ is Euclidean area. We maximize
\[
 k(f'(0),f'(0))+P_q(f)-\delta A_h(f)
\]
over finite-area discs with a fixed center, where $h$ is a K\"ahler
metric and $\delta>0$. The negative area term gives compactness. Comparing
a maximizing disc with its dilations then gives a linear boundary-area
tail bound. This improves its regularity to continuity on the closed
disc and $W^{1,p}$ for some $p>2$.

The main analytic work is to make boundary-normalized variations at
such an extremum legitimate. A weighted Hardy-space construction, in the tradition of matrix
spectral factorization \cite{WienerMasani1957,HelsonLowdenslager1958}, gives
a holomorphic frame $B=(b_1,\ldots,b_n)$ of the pulled-back tangent bundle with
$B^*hB=I$ on the boundary. We retain enough Sobolev regularity to realize
the sections $zb_j(z)$ as derivatives of actual center-fixed
holomorphic families $F_t$ with $F_0=f$. Let $\cL$ denote the sum of
the complex second derivatives in these $n$ parameter directions,
evaluated at $t=0$. The trace of the
complex Hessian of the area functional is then
\[
 \cL A_h(F_t)=n-A_{R_h}(f),
 \qquad R_h=-\ddc\log\det h,
 \qquad dd^c=i\partial\bar\partial.
\]
Its proof uses an explicit weak boundary limit; it does not assume that
the maximizing disc extends holomorphically across the circle. The
frame and variation constructions apply to boundary-regular
holomorphic discs independently of the global hyperbolicity
hypothesis, which is used to produce those discs.

Two choices of parameters exploit this identity:
\begin{center}
\begin{tabular}{@{}ll@{}}
\toprule
Parameters & Result \\
\midrule
$h=k$, $q=R_k$, $\delta\downarrow0$ & Uniform upper bound for $P_{R_k}$ \\
$q=-h$, $\delta=1$, $R_h\ge-h$ & Uniform lower bound for $h^n/k^n$ \\
\bottomrule
\end{tabular}
\end{center}
The first bound holds on all finite-area discs. The
Poletsky--Rosay disc envelope \cite{Poletsky1991,Rosay2003,DrinovecForstneric2012}
on $\det T^{1,0}X=K_X^*$ with its zero section removed then gives
bounded local plurisubharmonic weights for $K_X$. We include a direct
smoothing argument proving that $K_X$ is nef. Second, $q=-h$ gives the
pointwise estimate
\begin{equation}\label{eq:intro-volume}
 h^n\ge e^{-D}\left(\frac{n}{2n+D}\right)^n k^n
 \quad\text{whenever }R_h\ge-h.
\end{equation}
This estimate is uniform over the indicated metrics, even though the
auxiliary maximizing discs and their regularity constants need not be.

\subsection{From the estimates to ampleness}

Following the Monge--Amp\`ere and volume strategy used by Wu and Yau
\cite{WuYau2016}, the negative-sign Aubin--Yau equation produces K\"ahler metrics $h_t$, for $t>0$,
in the classes $2\pi c_1(K_X)+t[k]$ with $R_{h_t}=-h_t+t k$.
Equation~\eqref{eq:intro-volume} gives a positive lower bound for their
volumes. Passing to the limit in cohomology as $t\downarrow0$ proves
$\int_X c_1(K_X)^n>0$.

The remaining steps use established results with their usual
hypotheses. Holomorphic Morse inequalities make the nef canonical
bundle \emph{big}, meaning that its sections have maximal,
$n$-dimensional asymptotic growth. The resulting sections make $X$
Moishezon, that is, its meromorphic function field has transcendence
degree $n$. A compact K\"ahler Moishezon manifold is projective.
Finally, the absence of rational curves and the log cone theorem
upgrade bigness to ampleness. The final implication is recorded in
\cite[Lemma~2.1]{DiverioTrapani2019}; see also
\cite[Lemma~5.1]{Diverio2020}. All the new estimates are proved
below; the standard existence, envelope, and positivity theorems are
stated at the points where they are used.

\subsection{Organization}

Section~\ref{sec:preliminaries} fixes the curvature, Green-function,
and disc-functional conventions. Section~\ref{sec:extremal-discs}
constructs maximizing discs and proves their boundary regularity.
Section~\ref{var:section} builds the boundary-unitary frames and
centered holomorphic variations, and Section~\ref{sec:hessian}
establishes the determinant and Hessian identities. Their first use,
in Section~\ref{sec:nefness}, produces bounded canonical potentials
and proves nefness; their second use, in Section~\ref{sec:volume},
gives the uniform volume estimate. Section~\ref{sec:positive-intersection}
turns that estimate into positive top self-intersection, and
Section~\ref{sec:ampleness} completes the passage to bigness,
projectivity, and ampleness. Section~\ref{sec:consequences} applies the
result to pluricanonical sections and semialgebraic universal covers.
Appendix~\ref{app:smoothing} gives the bounded-potential smoothing
argument used in Section~\ref{sec:nefness}.

\section{Conventions and disc functionals}\label{sec:preliminaries}

Throughout, $X$ is a compact connected K\"ahler manifold of complex
dimension $n\ge1$, with no nonconstant entire curves, and $k$ is a fixed
K\"ahler metric. This section fixes the signs and normalizations used
in the disc estimates. We write $dd^c=i\partial\bar\partial$ and
identify a Hermitian tensor with its real
$(1,1)$-form so that the tensor associated with a real-valued potential
$\phi$ has quadratic form
\[
 \sum_{i,j}\partial_i\partial_{\bar j}\phi\,
 v^i\overline{v^j},
\]
and associated form $\ddc\phi$. In matrix formulas, vectors are columns
and their squared norm is written $v^*hv$. In particular, if the columns
of $B$ are tangent vectors, their Gram matrix is $B^*hB$.

For a K\"ahler form $h$, set
\begin{equation}\label{eq:ricci-convention}
 R_h=-\ddc\log\det h.
\end{equation}
This globally defined form represents $2\pi c_1(T^{1,0}X)$; hence
$[-R_h]=2\pi c_1(K_X)$. Inequalities of real $(1,1)$-forms are understood
as inequalities of their Hermitian tensors. The form $h^n$ is used
without a factor $n!$; determinant ratios are consequently ratios of
these volume forms.

Write $\D=\{z\in\C:|z|<1\}$ and let $\dd S$ denote Euclidean area.
For a holomorphic map $f:\D\to X$, put
\[
 f'(z)=\dd f_z(\partial_z),\qquad d(f)=k(f'(0),f'(0)).
\]
The disc has \emph{finite area} if $\int_\D k(f',f')\,\dd S<\infty$.
Since $X$ is compact, this condition is independent of the smooth
Hermitian metric used to define it. For a smooth closed real
$(1,1)$-form $q$ and a finite-area disc, define
\begin{equation}\label{eq:disc-functionals}
 A_q(f)=\frac1\pi\int_\D q(f',f')\,\dd S,
 \qquad
 P_q(f)=\frac2\pi\int_\D\log\frac1{|z|}\,q(f',f')\,\dd S.
\end{equation}
Both are finite: the weight is bounded away from the center, and the
integrand apart from the weight is smooth near the center. We write
$A_{q,r}(f)$ for the same area integral restricted to $|z|<r$.

For a scalar function on a circle, use the normalized mean
\[
 \langle u\rangle_r=\frac1{2\pi}\int_0^{2\pi}u(re^{i\theta})\,\dd\theta.
\]
The following normalization will be useful repeatedly.

\begin{lemma}[Green identities]\label{pre:green}
If $u$ is smooth on $\D$, then, for $0<r<1$,
\begin{align}
 \langle u\rangle_r-u(0)
 &=\frac2\pi\int_{|z|<r}\log\frac r{|z|}\,
       \partial_z\partial_{\bar z}u\,\dd S,
       \label{eq:green-poisson}\\
 \frac1\pi\int_{|z|<r}\partial_z\partial_{\bar z}u\,\dd S
 &=\frac r2\langle\partial_r u\rangle_r.
       \label{eq:green-area}
\end{align}
If in addition $u$ is continuous on $\overline\D$ and
$\partial_z\partial_{\bar z}u\in L^a(\D)$ for some $a>1$, then
Equation~\eqref{eq:green-poisson} passes to $r=1$.
\end{lemma}

\begin{proof}
The identities are the ordinary planar Green identities with
$\Delta=4\partial_z\partial_{\bar z}$. For the limiting assertion,
$\log(1/|z|)$ belongs to every finite $L^b(\D)$, so H\"older's
inequality and dominated convergence apply to the right-hand side.
Continuity gives convergence of the circle means.
\end{proof}

Sobolev regularity of a continuous map on $\overline\D$ is understood
in finitely many target coordinate charts. The same convention applies
to sections of a pulled-back bundle. When $p>2$, the embedding
$W^{1,p}(\D)\hookrightarrow C^0(\overline\D)$, its local versions, and
the multiplication estimate make $W^{1,p}$ an algebra. Smooth
functions of finitely many Sobolev coefficients are again Sobolev
when their values stay in a compact subset of their domain. We use the
corresponding smooth composition maps between these Banach spaces.
For fiberwise holomorphic functions with uniform radii and Sobolev
coefficient bounds, composition is holomorphic between the complex
Banach spaces in question; the precise application is given in the
variation construction.

The following table collects the notation used throughout the disc
argument. The area and Poisson functionals use the fixed unit disc,
with $A_{q,r}$ denoting the explicitly indicated truncation.
\begin{center}
\begin{tabular}{@{}ll@{}}
\toprule
Notation & Meaning \\
\midrule
$k$ & Fixed background K\"ahler metric on $X$ \\
$h$ & Auxiliary K\"ahler metric \\
$R_h$ & $-dd^c\log\det h$, representing $-2\pi c_1(K_X)$ \\
$d(f)$ & Squared central derivative norm $k(f'(0),f'(0))$ \\
$A_q(f)$ & $\pi^{-1}\int_\D q(f',f')\,\dd S$ \\
$A_{q,r}(f)$ & The same area integral over $|z|<r$ \\
$P_q(f)$ & $2\pi^{-1}\int_\D\log(1/|z|)q(f',f')\,\dd S$ \\
$\langle u\rangle_r$ & Normalized mean of $u$ on $|z|=r$ \\
\bottomrule
\end{tabular}
\end{center}

\section{Extremal discs and boundary regularity}
\label{sec:extremal-discs}

Our goal is to maximize the disc functional and obtain enough boundary
regularity to vary the whole disc. Hyperbolicity first gives compactness
on interior subdiscs; an area penalty will control the boundary tail.
Throughout this section, $k$ is the fixed K\"ahler metric on $X$.
We write $|f'(z)|_k^2=k(f'(z),f'(z))$ and use the path distance of
the associated Riemannian metric when discussing uniform convergence.
The fixed numerical factor relating real and complex tangent norms will
be absorbed into constants.

\begin{lemma}[Derivative bound and normality]
\label{ex:brody}
There is a constant $D>0$, depending only on $(X,k)$, such that every
holomorphic map $f\colon\D\to X$ satisfies
\[
 d(f)=|f'(0)|_k^2\leq D.
\]
Consequently,
\begin{equation}
 |f'(z)|_k\leq\frac{\sqrt D}{1-|z|},\qquad z\in\D.
 \label{eq:brody-derivative}
\end{equation}
Every sequence of such maps has a subsequence converging uniformly on
compact subsets of $\D$ to a holomorphic map into $X$. In local target
coordinates, all derivatives converge uniformly on smaller compact
subsets as well.
\end{lemma}

\begin{proof}
We give Brody's rescaling argument in the normalization needed here
\cite{Brody1978}.
Suppose first that there are maps $f_j\colon\D\to X$ with
$|f_j'(0)|_k\to\infty$. Choose $z_j$ at which the continuous function
\[
 z\longmapsto (1/2-|z|)|f_j'(z)|_k
\]
attains its maximum on $\{|z|\leq1/2\}$. For all sufficiently large $j$
this maximum is positive, so $|z_j|<1/2$. Put
\[
 \lambda_j=|f_j'(z_j)|_k,
 \qquad R_j=(1/2-|z_j|)\lambda_j,
 \qquad g_j(w)=f_j(z_j+w/\lambda_j).
\]
The map $g_j$ is defined on $\{|w|<R_j\}$, and
\[
 R_j\geq\tfrac12|f_j'(0)|_k\longrightarrow\infty,
 \qquad |g_j'(0)|_k=1.
\]
If $|w|<R_j/2$, then
\[
 1/2-|z_j+w/\lambda_j|
 \geq\tfrac12(1/2-|z_j|).
\]
The maximizing property of $z_j$ therefore gives
$|g_j'(w)|_k\leq2$ on this smaller disc.

For every fixed radius these maps are consequently equicontinuous and
take values in the compact metric space $X$. The Arzel\`a--Ascoli theorem
and a diagonal subsequence give a locally uniform limit
$g\colon\C\to X$. To verify holomorphy, choose a point of the source and
a target coordinate chart containing its image under $g$. On a
sufficiently small source neighborhood, the limit and all sufficiently
late maps take values in that chart. Their coordinate representations
converge uniformly on smaller neighborhoods. The usual theorem on
uniform limits of holomorphic functions, followed by the Cauchy
integral formula for derivatives, shows that $g$ is holomorphic and
$|g'(0)|_k=1$. This contradicts the hypothesis on entire maps. The
central derivative bound follows.

For $z\in\D$, apply this bound to
$w\mapsto f(z+(1-|z|)w)$. Its central derivative is
$(1-|z|)f'(z)$, proving Equation~\eqref{eq:brody-derivative}. The estimate gives
equicontinuity on every compact source disc. A second application of
Arzel\`a--Ascoli and the same coordinate argument proves normality and
convergence of derivatives. Increasing $D$ if necessary makes it
strictly positive.
\end{proof}

We next fix a point $x\in X$, a smooth K\"ahler form $h$, a smooth
closed real $(1,1)$-form $q$, and a number $\delta>0$. All constants in
the next two results may depend on these choices. Finite area can be
measured using either $h$ or $k$, since the two metrics are uniformly
comparable on $X$. We consider
\begin{equation}
 J_{h,q,\delta}(f)=d(f)+P_q(f)-\delta A_h(f)
 \label{eq:disc-functional}
\end{equation}
on the finite-area holomorphic discs with $f(0)=x$.

\begin{proposition}[Existence of a maximizing disc]
\label{ex:maximizer}
The functional in Equation~\eqref{eq:disc-functional} has finite
supremum, attained by a finite-area holomorphic disc centered at $x$.
\end{proposition}

\begin{proof}
Choose $C_q\geq0$ such that
\[
 |q(v,v)|\leq C_q h(v,v)
\]
for every tangent vector $v$. Select $r_0\in[1/2,1)$ so close to $1$
that $2C_q\log(1/r_0)\leq\delta/2$. For a disc $f$ write
\[
 a_f(z)=h(f'(z),f'(z)),\qquad b_f(z)=q(f'(z),f'(z)).
\]
On the outer annulus $r_0<|z|<1$, the function
\[
 T_f(z)=\delta a_f(z)-2\log(1/|z|)b_f(z)
\]
satisfies
\begin{equation}
 T_f(z)\geq\frac\delta2 a_f(z)\geq0.
 \label{eq:negative-tail-density}
\end{equation}
On the inner disc, Lemma~\ref{ex:brody} and comparison of $h$ with $k$
give a uniform bound for $a_f$ and $|b_f|$. Since
$\log(1/|z|)$ is integrable there, there is a constant $K$ independent
of $f$ such that
\[
 \left|\frac1\pi\int_{|z|<r_0}
 \bigl(2\log(1/|z|)b_f(z)-\delta a_f(z)\bigr)\,dS\right|
 \leq K.
\]
Consequently,
\begin{equation}
 J_{h,q,\delta}(f)
 \leq D+K-\frac\delta{2\pi}\int_{r_0<|z|<1}a_f\,dS.
 \label{eq:maximizing-area-bound}
\end{equation}
In particular, the supremum is finite. It is nonnegative because the
constant disc has value zero.

Choose a maximizing sequence $(f_j)$ with $J_{h,q,\delta}(f_j)\geq-1$.
Equation~\eqref{eq:maximizing-area-bound} bounds its outer-annulus
areas uniformly, and the compact-interior derivative estimate bounds
its inner areas. Thus
\[
 \sup_j A_h(f_j)<\infty.
\]
By Lemma~\ref{ex:brody}, after passing to a subsequence the maps converge
locally uniformly to a holomorphic disc $f$ centered at $x$. Derivative
convergence on compact subsets gives pointwise convergence of $a_{f_j}$
and $b_{f_j}$. Fatou's lemma gives
$A_h(f)\leq\liminf_j A_h(f_j)<\infty$.

We verify the required upper semicontinuity of $J$. Its central term
converges by derivative convergence. Its inner-disc integral converges
by dominated convergence, using the uniform bound above together with
the integrable logarithmic weight. On the outer annulus, the
nonnegative functions $T_{f_j}$ converge pointwise to $T_f$, so Fatou's
lemma yields
\[
 -\frac1\pi\int_{r_0<|z|<1}T_f\,dS
 \geq\limsup_j
 \left(-\frac1\pi\int_{r_0<|z|<1}T_{f_j}\,dS\right).
\]
All these integrals are finite. Indeed, finite area controls the
unweighted integrals and the logarithmic weight is bounded on the
outer annulus. Combining the three terms proves
\[
 J_{h,q,\delta}(f)\geq\limsup_jJ_{h,q,\delta}(f_j),
\]
which proves attainment.
\end{proof}

Recall that $A_{q,s}(f)$ is the area integral restricted to $|z|<s$.
The form $q$ need not be positive, so this quantity can have either
sign. Nevertheless, $|q|\leq C_qh$ implies that $A_{q,s}(f)$ is bounded
in $s$ and tends to $A_q(f)$ as $s$ tends to $1$.
Maximality now improves the finite area of the disc to a quantitative
tail estimate. This is the additional control needed at the boundary.

\begin{lemma}[Dilation and area tails]
\label{ex:dilation}
If $f$ is a maximizer in Proposition~\ref{ex:maximizer}, then, for
$0<r<1$,
\begin{equation}
 \begin{split}
 0\leq\delta\bigl(A_h(f)-A_{h,r}(f)\bigr)
 &\leq(1-r^2)d(f)
       +2\int_r^1 A_{q,s}(f)\,\frac{ds}{s}.
 \end{split}
 \label{eq:tail}
\end{equation}
In particular,
\begin{equation}
 A_q(f)\geq-d(f),\qquad
 A_k(f)-A_{k,r}(f)\leq C_f(1-r)\quad(1/2\leq r<1)
 \label{eq:extremal-area-tail}
\end{equation}
for some finite constant $C_f$.
\end{lemma}

\begin{proof}
Let $f_r(z)=f(rz)$. A change of variables gives
\[
 d(f_r)=r^2d(f),\qquad
 A_h(f_r)=A_{h,r}(f),\qquad
 A_{q,s}(f_r)=A_{q,rs}(f).
\]
The identity $\log(1/|z|)=\int_{|z|}^1ds/s$ and Fubini's theorem give
\[
 P_q(f)=2\int_0^1 A_{q,s}(f)\,\frac{ds}{s}.
\]
Fubini is valid also for this signed integrand: near the origin the
coordinate derivatives of $f$ are bounded, while away from the
origin the logarithmic weight is bounded and $|q(f',f')|$ is controlled
by the finite $h$-area. Applying the same identity to $f_r$ gives
\begin{equation}
 P_q(f)-P_q(f_r)
   =2\int_r^1 A_{q,s}(f)\,\frac{ds}{s}.
 \label{eq:dilation-poisson}
\end{equation}
Since $f_r$ has the same center and finite area, maximality gives
$J_{h,q,\delta}(f)\geq J_{h,q,\delta}(f_r)$. Substituting the three
identities proves Equation~\eqref{eq:tail}; its leftmost inequality uses the
positivity of $h$.

Divide the rightmost expression in Equation~\eqref{eq:tail} by $1-r$. As
$r\uparrow1$, it tends to $2d(f)+2A_q(f)$. It is nonnegative for every
$r$, and hence $A_q(f)\geq-d(f)$. For $r\geq1/2$, the original
right-hand side of Equation~\eqref{eq:tail} is at most
\[
 2d(f)(1-r)+4C_qA_h(f)(1-r).
\]
This bounds the $h$-area tail by a constant times $1-r$. Uniform
comparison of $k$ and $h$ then gives the second assertion in
Equation~\eqref{eq:extremal-area-tail}.
\end{proof}

The tail estimate supplies the last step of this section: continuity
on the closed disc and a Sobolev exponent strictly above two. These
are precisely the hypotheses used to construct the frames and
variations in Section~\ref{var:section}.
\begin{proposition}[Boundary regularity of maximizing discs]
\label{ex:regularity}
Every disc maximizing the functional in Equation~\eqref{eq:disc-functional}
extends continuously
to $\overline\D$ and belongs to $W^{1,p}(\D,X)$ for every $2<p<4$.
In particular, it has continuous $W^{1,3}$ regularity up to the
boundary. More precisely, there is a constant $C_f$ such that
\begin{equation}
 |f'(z)|_k\leq C_f(1-|z|)^{-1/2}
 \label{eq:boundary-derivative}
\end{equation}
for $z$ sufficiently close to $\partial\D$.
\end{proposition}

\begin{proof}
We first make precise the uniform choice of target charts for small
source discs. Choose finitely many nested coordinate neighborhoods
\[
 V_\alpha\Subset W_\alpha\Subset U_\alpha
\]
such that the $V_\alpha$ cover $X$ and each $U_\alpha$ is a holomorphic
coordinate chart. By compactness, there is a number $\rho>0$ such that
if $y\in V_\alpha$, its metric ball of radius $\rho$ is contained in
$W_\alpha$. On each $\overline{W_\alpha}$, the metric $k$ and the
Euclidean coordinate metric are uniformly comparable; one comparison
constant suffices for this finite collection.

For $z\in\D$ put $t=1-|z|$. Let $0<c<1/4$, to be fixed independently
of $z$ and $f$. For $w\in B(z,ct)$, the straight segment from $z$ to
$w$ stays at source distance at least $(1-c)t$ from $\partial\D$.
Equation~\eqref{eq:brody-derivative} therefore bounds the length of
its image by
\[
 C\sqrt D\,\frac{c}{1-c},
\]
where $C$ accounts only for the convention relating real and complex
tangent norms. Choose $c$ small enough that this bound is less than
$\rho$. Selecting $\alpha$ with $f(z)\in V_\alpha$, we obtain
\[
 f\bigl(B(z,ct)\bigr)\subset W_\alpha.
\]

In these coordinates the derivative of $f$ is a vector of holomorphic
functions. Its squared Euclidean norm is subharmonic. The mean-value
inequality and the uniform metric comparisons give
\[
 |f'(z)|_k^2
 \leq\frac{C_1}{t^2}
       \int_{B(z,ct)}|f'(w)|_k^2\,dS(w).
\]
Every point of $B(z,ct)$ has modulus greater than
$1-(1+c)t$. For $t$ sufficiently small, the area-tail estimate in
Equation~\eqref{eq:extremal-area-tail} consequently gives
\[
 \int_{B(z,ct)}|f'(w)|_k^2\,dS(w)
 \leq\pi\bigl(A_k(f)-A_{k,1-(1+c)t}(f)\bigr)
 \leq C_2t.
\]
This proves Equation~\eqref{eq:boundary-derivative}.

For $r<s<1$ sufficiently close to $1$ and $\zeta\in\partial\D$,
integration along the radial segment gives
\[
 \operatorname{dist}_k\bigl(f(r\zeta),f(s\zeta)\bigr)
 \leq C_3\int_r^s(1-u)^{-1/2}\,du
 \leq2C_3\sqrt{1-r}.
\]
The estimate is uniform in $\zeta$. Since $X$ is complete, each radial
path has a limit, and the continuous maps
$\zeta\mapsto f(r\zeta)$ converge uniformly as $r\uparrow1$. Their
limit is therefore continuous on $\partial\D$. The same uniform
estimate, together with continuity of this boundary map, proves that
the resulting extension of $f$ is continuous on all of $\overline\D$.

It remains to improve the integrability of the derivative. For all
sufficiently large integers $j$, let
\[
 E_j=\{z\in\D:2^{-j-1}<1-|z|<2^{-j}\}.
\]
The area-tail estimate and Equation~\eqref{eq:boundary-derivative} imply
\[
 \int_{E_j}|f'|_k^2\,dS\leq C_4\,2^{-j},
 \qquad
 \sup_{E_j}|f'|_k\leq C_5\,2^{j/2}.
\]
Thus, for $2<p<4$,
\[
 \begin{split}
 \int_{E_j}|f'|_k^p\,dS
 &\leq\bigl(\sup_{E_j}|f'|_k\bigr)^{p-2}
              \int_{E_j}|f'|_k^2\,dS\\
 &\leq C_p\,2^{-j(4-p)/2}.
 \end{split}
\]
This is summable in $j$. On the remaining compact interior disc the
derivative is bounded, so $|f'|_k\in L^p(\D)$.

For completeness, the asserted Sobolev regularity is understood in
local target coordinates. Continuity on $\overline\D$ allows a finite
cover of the closed source disc by small neighborhoods whose images
lie in relatively compact target charts. In each such chart, the
coordinate functions are bounded, are smooth in the source interior,
and have first derivatives in $L^p$ by the estimate just proved and
metric comparison. These classical derivatives are their weak
derivatives, as is seen by testing against compactly supported smooth
functions in the source interior. Hence the coordinate functions lie
in $W^{1,p}$ up to the boundary, with the continuous boundary values
already constructed. This proves the proposition.
\end{proof}

\begin{remark}
The constants in Proposition~\ref{ex:regularity} need not be uniform
in $h$, $q$, or $\delta$. What will be used is the regularity of each
individual maximizing disc, to construct variations with unrestricted
boundary values. The constant $D$ in Lemma~\ref{ex:brody}, in contrast,
depends only on the fixed metric $k$ and the compact manifold $X$.
\end{remark}

\section{Boundary frames and centered holomorphic variations}
\label{var:section}

We establish the two analytic facts needed to vary the maximizing disc.
Throughout this section, let
\[
 f\colon\overline\D\longrightarrow X
\]
be continuous, holomorphic on $\D$, and of class $W^{1,p}$ for some
$p>2$, with the coordinate convention of Section~\ref{sec:preliminaries}. The
maximizers supplied by Proposition~\ref{ex:regularity} satisfy these
assumptions.  No extension of $f$ across $\partial\D$ is assumed.

We use two elementary consequences of $p>2$: the embedding
$W^{1,p}(\D)\subset C^0(\overline\D)$ and the fact that
$W^{1,p}(\D)$ is a Banach algebra under pointwise multiplication.
In particular, a continuous, pointwise invertible $W^{1,p}$ matrix on
$\overline\D$ has a $W^{1,p}$ inverse.

\subsection{A holomorphic frame regular at the boundary}

We first trivialize the pulled-back tangent bundle while retaining
the boundary Sobolev regularity of $f$. We will then change this
frame so that its boundary values are orthonormal for $h$.
\begin{lemma}\label{var:initial-frame}
The bundle $f^*T^{1,0}X$ admits a holomorphic frame $e$ on $\D$ which
extends to a continuous, nonsingular $W^{1,p}$ frame on $\overline\D$.
\end{lemma}

\begin{proof}
The continuous pullback bundle over $\overline\D$ is topologically
trivial.  Its coordinate transition matrices are $W^{1,p}$: they are
smooth functions of $f$ with bounded derivatives on the compact subsets
of the target charts in use.  Start with a continuous frame, approximate
its coefficients uniformly in finitely many local frames, and combine
the approximations with a smooth partition of unity on the source.
This gives a $W^{1,p}$ frame $v$ uniformly close to the original one;
sufficiently close approximation preserves its pointwise independence.

In the frame $v$, the Dolbeault operator on the open disc has the form
\[
 \bar\partial+a,
 \qquad a\in L^p(\D,\operatorname{Mat}_n(\C)),
\]
where we identify a $(0,1)$-form with its coefficient.  Extend $a$ by
zero to a disc $\D_R=\{|z|<R\}$ with $R>1$.  On a sufficiently small
disc $U\subset\D_R$ of radius $r$, solve
\begin{equation}\label{var:local-gauge}
 \bar\partial U_0=-aU_0.
\end{equation}
Here and below the same notation for a matrix equation is understood
entrywise.  For completeness, the solid Cauchy transform
\[
 T_Ug(z)=\frac1\pi\int_U\frac{g(\zeta)}{z-\zeta}\,\dd S(\zeta)
\]
satisfies $\bar\partial T_Ug=g$ weakly in $U$.  H\"older's inequality
for its kernel and the $L^p$ estimate for the planar Cauchy singular
integral~\cite[Section~4.1 of the authors' preprint]{SukhovTumanov2016} give
\begin{equation}\label{var:cauchy-estimates}
 \|T_Ug\|_{L^\infty(U)}
       \le C_p r^{1-2/p}\|g\|_{L^p(U)},
 \qquad
 \|T_Ug\|_{W^{1,p}(U)}\le C_{p,r}\|g\|_{L^p(U)}.
\end{equation}
Thus $U_0=I-T_U(aU_0)$ is solved by contraction in the sup norm when
$r$ is small.  The radius can also be chosen so that the solution stays
within distance $1/2$ of $I$.  It is therefore invertible and, by
Equation~\eqref{var:cauchy-estimates}, belongs to $W^{1,p}$.

For any two such solutions $U_i,U_j$, the transition matrix
$U_i^{-1}U_j$ satisfies $\bar\partial(U_i^{-1}U_j)=0$ weakly on the
overlap.  It is consequently holomorphic there.  These transitions
define a holomorphic vector bundle on $\D_R$.  A holomorphic vector
bundle on a disc is holomorphically trivial, by the Oka--Grauert
principle~\cite{Grauert1958,ForstnericPrezelj2000}.  Choose a global holomorphic frame of
this bundle.  Relative to its local frames $U_i$, its coefficients are
holomorphic; relative to the original trivial bundle they are therefore
$W^{1,p}$ on compact subsets of $\D_R$.  In particular the resulting
matrix $U$ is continuous, invertible, and $W^{1,p}$ on $\overline\D$,
and satisfies $\bar\partial U=-aU$ on $\D$.  The frame $e=vU$ has all
the required properties.
\end{proof}

\begin{proposition}[Boundary-unitary frame]\label{var:frame}
Let $h$ be a smooth Hermitian metric on $X$.  There exist a number
$p_0$ with $2<p_0\le p$ and a holomorphic frame
$B=(b_1,\ldots,b_n)$ of $f^*T^{1,0}X$ which is continuous,
nonsingular, and $W^{1,p_0}$ on $\overline\D$, such that
\begin{equation}\label{var:boundary-unitary}
 B(z)^*h(f(z))B(z)=I,
 \qquad z\in\partial\D.
\end{equation}
\end{proposition}

The factorization step belongs to the matrix spectral-factorization
theory of Wiener--Masani and Helson--Lowdenslager
\cite{WienerMasani1957,HelsonLowdenslager1958}. We give the weighted
Hardy-space construction and establish the boundary Sobolev regularity
required for our variations.

\begin{proof}
Choose $e$ by Lemma~\ref{var:initial-frame} and set
\[
 M(z)=e(z)^*h(f(z))e(z).
\]
This is a continuous, positive definite $W^{1,p}$ matrix on the closed
disc.  In particular, there are constants $0<m\le M_0<\infty$ such
that
\begin{equation}\label{var:metric-bounds}
 mI\le M(z)\le M_0I
 \quad\hbox{on }\overline\D.
\end{equation}

\paragraph{The Hardy-space factorization.}
Write $H^2=H^2(\partial\D,\C^n)$ for the vector Hardy space with
nonnegative Fourier modes, and put $\dd\sigma=\dd\theta/(2\pi)$.
Equip the same vector space with the equivalent Hilbert norm
\[
 \|u\|_M^2=\int_{\partial\D}u^*Mu\,\dd\sigma.
\]
Multiplication by $z$ is an isometry, and its closed image $zH^2$ has
codimension $n$.  Let $s_1,\ldots,s_n$ be an orthonormal basis of
$W=H^2\ominus_M zH^2$ and form the matrix $S=(s_1,\ldots,s_n)$.
The vectors $z^j s_a$, with $j\ge0$ and $1\le a\le n$, are an
orthonormal basis.  Indeed, iterating
$H^2=W\oplus_M zH^2$ shows that the orthogonal complement of their
span is contained in every $z^jH^2$; the intersection of these spaces
is zero.

Their orthogonality identifies all Fourier coefficients of the
$L^1$ matrix $S^*MS$, and gives
\begin{equation}\label{var:hardy-factor}
 S^*MS=I\quad\hbox{almost everywhere on }\partial\D.
\end{equation}
Equation~\eqref{var:metric-bounds} then bounds the boundary values of
$S$ in $L^\infty$.  Its holomorphic Hardy extension is bounded on
$\D$, as follows entrywise from the Poisson integral.  It is invertible
at every point $\zeta\in\D$: evaluation at $\zeta$ is a continuous
surjection from $H^2$ onto $\C^n$, whereas the evaluations of all
$z^j s_a$ lie in the range of $S(\zeta)$.  Their dense span forces
that range to be $\C^n$.

\paragraph{Regularity of the boundary factor.}
The factor constructed so far has only almost-everywhere boundary
values. We now obtain continuity and a Sobolev exponent above two,
which will also give nonsingularity at every boundary point.
We first record the fractional regularity supplied by $M\in W^{1,p}$.
If $M_\theta(e^{i\alpha})=M(e^{i(\alpha+\theta)})$, then
\begin{equation}\label{var:trace-difference}
 \|M_\theta-M\|_{L^p(\partial\D)}
       \le C|\theta|^{1-1/p},\qquad |\theta|\le\tfrac14.
\end{equation}
One can see this directly, without a trace theorem.  Set
$t=|\theta|$ and choose $\rho\in[1-2t,1-t]$ such that the angular
$L^p$ norm of $\partial_\alpha M(\rho e^{i\alpha})$ is at most
$Ct^{-1/p}\|M\|_{W^{1,p}}$.  Radial integration and H\"older's
inequality bound
$\|M|_{\partial\D}-M(\rho\,\cdot)\|_{L^p}$ by
$Ct^{1-1/p}\|M\|_{W^{1,p}}$.  At radius $\rho$, angular integration
gives the same bound for translation through $\theta$.  The triangle
inequality proves Equation~\eqref{var:trace-difference}; the argument
for Sobolev functions follows by approximation, with the continuous
boundary values identifying the trace.

Let $H^s(S^1)$ now denote the fractional Sobolev space on the circle,
defined by square summability of Fourier coefficients with weight
$1+|l|^{2s}$. This differs from the Hardy-space notation $H^2$ above.
It follows that $M|_{\partial\D}\in H^s(S^1)$ for every
\begin{equation}\label{var:s-choice}
 \frac12<s<1-\frac1p.
\end{equation}
Indeed, on a circle of finite measure the $L^p$ bound controls the
$L^2$ bound, and integration of the squared rotation differences
against $|\theta|^{-1-2s}\,\dd\theta$ is finite.  Parseval's identity
identifies this condition with
$\sum_{l\in\mathbb Z}(1+|l|^{2s})|M_l|^2<\infty$.

Let $\Pi$ denote the unweighted orthogonal Hardy projection.  The
Toeplitz operator
\[
 T_Mv=\Pi(Mv),\qquad T_M\colon H^2\longrightarrow H^2,
\]
is bounded, self-adjoint, and satisfies
$\langle T_Mv,v\rangle\ge m\|v\|_2^2$.  Consequently it has a
bounded inverse: its range is closed by the lower bound and dense
because its orthogonal complement is the kernel of its adjoint.
For a column $v$ of $S$, membership in $W$ says precisely that
$T_Mv$ is a constant vector.  Rotating this identity and subtracting
it from the original one yields
\[
 T_M(v_\theta-v)=\Pi\bigl((M-M_\theta)v_\theta\bigr).
\]
The boundedness of $v$ and $T_M^{-1}$ gives
\begin{equation}\label{var:factor-difference}
 \|v_\theta-v\|_2
 \le C\|(M_\theta-M)v_\theta\|_2
 \le C\|M_\theta-M\|_2.
\end{equation}
The same rotation characterization therefore gives $v\in H^s$.

Write $v(z)=\sum_{l\ge0}v_lz^l$.  Since $s>1/2$, Cauchy--Schwarz
shows that this series converges absolutely and uniformly on
$\overline\D$.  Thus $v$ is continuous there.  More precisely, for
$1/2<r<1$, Parseval's identity and Cauchy--Schwarz give
\begin{align}
 \|v'(r\,\cdot)\|_2^2
 &\le C(1-r)^{2s-2}
       \sum_{l\ge1}l^{2s}|v_l|^2,
       \label{var:factor-l2}\\
 \|v'(r\,\cdot)\|_\infty
 &\le C(1-r)^{s-3/2}
       \left(\sum_{l\ge1}l^{2s}|v_l|^2\right)^{1/2}.
       \label{var:factor-sup}
\end{align}
For any $p_0>2$, interpolation of these two estimates bounds
$\|v'(r\,\cdot)\|_{p_0}^{p_0}$ by a constant times
\[
 (1-r)^{(s-3/2)(p_0-2)+2s-2}
       =(1-r)^{p_0(s-3/2)+1}.
\]
This is integrable in $r$ provided
\begin{equation}\label{var:p0-choice}
 2<p_0<\frac{2}{3/2-s}.
\end{equation}
The interval is nonempty because $s>1/2$, and we choose $p_0$ also
at most $p$.  It follows that $S\in W^{1,p_0}(\D)$.

By continuity, Equation~\eqref{var:hardy-factor} now holds at every
boundary point.  It implies that $S$ is invertible there as well as
in the interior.  Compactness gives uniform nonsingularity on
$\overline\D$.  Finally $B=eS$ is a continuous, nonsingular,
$W^{1,p_0}$ holomorphic frame and satisfies
Equation~\eqref{var:boundary-unitary}.  Notice that uniform
nonsingularity was obtained after boundary regularity; it was not
assumed in the Hardy-space argument.
\end{proof}

\subsection{Integrating the infinitesimal variations}

The boundary-unitary frame provides sections $zb_j$ vanishing at the
center. To use maximality, these sections must be derivatives of
admissible holomorphic discs. Sprays of analytic discs fixing an
exceptional set are standard in disc-envelope theory
\cite[Definition~2.1 and Lemma~2.2]{DrinovecForstneric2012}.
The following construction records the prescribed derivatives and
Sobolev parameter dependence required by the Hessian calculation.

\begin{proposition}[Centered holomorphic variations]\label{var:spray}
Let $B$ and $p_0>2$ be as in Proposition~\ref{var:frame}, and set
$x=f(0)$.  For some $\eta>0$, there is a map
\[
 F\colon\{t\in\C^n:|t|<\eta\}\times\overline\D\longrightarrow X
\]
which is continuous on its domain and jointly holomorphic on
$\{|t|<\eta\}\times\D$, and satisfies
\begin{equation}\label{var:spray-identities}
 F(0,z)=f(z),\qquad F(t,0)=x,\qquad
 \frac{\partial F}{\partial t_j}(0,z)=z b_j(z).
\end{equation}
In fixed local target coordinates covering the graph of $f$, the
maps $t\mapsto F(t,\cdot)$ are holomorphic with values in
$W^{1,p_0}$ on the corresponding source patches.  In particular
every $F(t,\cdot)$ is a finite-area holomorphic disc centered at $x$.
\end{proposition}

The proof first gives the graph of $f$ fiberwise holomorphic
coordinates with Sobolev dependence on the source. The resulting
coordinate maps need not be holomorphic in the source; we correct
that defect by a nonlinear $\bar\partial$ equation. Subtracting the
central value in its right inverse preserves the prescribed center.

\begin{proof}
\textbf{Fiber coordinates.}
Choose a frame $e$ from Lemma~\ref{var:initial-frame} and write
$B=eS$. The matrix $S=e^{-1}B$ is holomorphic on $\D$ and belongs
to $W^{1,p_0}$ on $\overline\D$. We first construct fiber coordinates near the graph of
$f$.  Choose a finite cover of $\overline\D$ by relatively open
sets $U_i$ and target charts
$\phi_i\colon V_i\longrightarrow\C^n$ so that $f(\overline U_i)$
is a compact subset of $V_i$, after shrinking the source sets as
necessary.  Let $\rho_i$ be a smooth partition of unity on the
closed disc subordinate to these sets.  In the $i$th chart put
\[
 A_i(z)=D\phi_i(f(z))e(z),\qquad
 C_i(z,y)=A_i(z)^{-1}\bigl(\phi_i(y)-\phi_i(f(z))\bigr).
\]
For $z\in U_i$, the matrix $A_i$ is holomorphic in the interior,
continuous and invertible up to the boundary, and of class $W^{1,p_0}$.
The map $C_i$ is holomorphic in $y$ and, for fixed admissible $y$,
holomorphic in the interior source variable $z$.  On a neighborhood
of the graph of $f$ define
\begin{equation}\label{var:forward-coordinates}
 \Phi(z,y)=\sum_i\rho_i(z)C_i(z,y).
\end{equation}
Each term is used only where its target chart is defined; the compact
support of $\rho_i$ inside the corresponding source patch makes
extension by zero harmless.  These neighborhoods can be chosen
uniformly about the graph.  More precisely, on sufficiently small
source patches all indices whose supports meet the patch have their
charts defined on one common target neighborhood of its image under
$f$.  This follows from finiteness of the cover and subordination of
the supports.  By construction,
\begin{equation}\label{var:coordinate-normalization}
 \Phi(z,f(z))=0,
 \qquad D_y\Phi(z,f(z))\,e(z)=I.
\end{equation}

\paragraph{Uniform inverses and their Sobolev dependence.}
Work on one of these smaller source patches, with fixed target chart
$\psi$, and put $c(z)=\psi(f(z))$ and
$e_\psi(z)=D\psi(f(z))e(z)$.  Normalize the forward map by
\[
 K_z(v)=\Phi\bigl(z,\psi^{-1}(c(z)+e_\psi(z)v)\bigr).
\]
Then $K_z(0)=0$ and $D_vK_z(0)=I$.  Continuity of all fiber
derivatives and compactness allow a common $R>0$ such that the maps
are defined on $|v|\le R$ and
\[
 \|D_vK_z(v)-I\|\le\tfrac14.
\]
For $|\xi|<R/2$, the equation
\[
 v=\xi-\bigl(K_z(v)-v\bigr)
\]
is a contraction of the closed ball of radius $R$ into itself.
Its solution $V(z,\xi)$ is continuous in $(z,\xi)$ and holomorphic
in $\xi$, as follows from its uniformly convergent contraction
iterations.  Define $E$ by
$\psi(E(z,\xi))=c(z)+e_\psi(z)V(z,\xi)$.  The local inverses
agree near $\xi=0$ by uniqueness and then throughout their common
fiber domain by holomorphy.  Taking the minimum of the finitely many
radii gives a common fiber neighborhood, with
\[
 \Phi(z,E(z,\xi))=\xi,
\]
and
\begin{equation}\label{var:inverse-normalization}
 E(z,0)=f(z),\qquad E_\xi(z,0)=e(z).
\end{equation}

We next obtain a Sobolev estimate uniform in the fiber parameter.
The chart formula for $K_z$ depends smoothly on the finite coefficient list
\[
 a(z)=\bigl(\rho_i(z),A_i(z)^{-1},\phi_i(f(z)),
                 c(z),e_\psi(z)\bigr)_i.
\]
All these coefficients are $W^{1,p_0}$ with compact range.  Fix a
smaller closed fiber polydisc inside $|\xi|<R/2$.  The strict bounds
in the contraction argument persist when the coefficient list varies
in a sufficiently small neighborhood of its compact range.  The
finite-dimensional implicit function theorem, with the coefficients
regarded as real parameters, therefore makes $V$ a smooth function of
that list and $\xi$, holomorphic in $\xi$.  After shrinking this
coefficient neighborhood, its first coefficient derivatives are
uniformly bounded on the fixed fiber polydisc.  The same is true for
the coordinate inverse
$E_\psi(z,\xi):=\psi(E(z,\xi))=c(z)+e_\psi(z)V(z,\xi)$.
The Sobolev chain rule now shows that $E_\psi(\cdot,\xi)$ is
$W^{1,p_0}$, with its weak source derivatives bounded by a single
$L^{p_0}$ function independent of $\xi$.  In particular, for some
$b\in L^{p_0}$ on the source patch,
\begin{equation}\label{var:inverse-sobolev-bound}
 |E_\psi(z,\xi)|\le C,
 \qquad |\nabla_zE_\psi(z,\xi)|\le C b(z)
 \quad\hbox{for almost every }z,
\end{equation}
uniformly for $\xi$ in that smaller polydisc.  One may take $b$ to
be a constant plus the sum of the absolute values of the first weak
derivatives of the finite coefficient list.  Derivatives with respect
to $\xi$ satisfy the same type of bound after another fixed shrink
of the polydisc, by Cauchy's formula.

\paragraph{Holomorphic substitution.}
These estimates also justify holomorphy of the substitution operator
\begin{equation}\label{var:inverse-substitution}
 u\longmapsto E_\psi(\cdot,u(\cdot))
\end{equation}
from a neighborhood of zero in $W^{1,p_0}(\D,\C^n)$ into the local
$W^{1,p_0}$ space.  Indeed, expand in the fiber variable:
\[
 E_\psi(z,\xi)=\sum_{\alpha\in\mathbb N^n}E_\alpha(z)\xi^\alpha.
\]
Cauchy's formula on a fixed fiber polydisc, with weak $z$
differentiation under its contour integral, gives
\[
 \|E_\alpha\|_\infty+\|\nabla E_\alpha\|_{p_0}
       \le Cr_1^{-|\alpha|}
\]
for some $r_1>0$, by Equation~\eqref{var:inverse-sobolev-bound}.
The algebra norm $\|a\|_\infty+\|\nabla a\|_{p_0}$ is submultiplicative.
Thus $\sum_\alpha E_\alpha u^\alpha$ converges normally in the
local $W^{1,p_0}$ space when the global $W^{1,p_0}$ norm of $u$ is small;
the number of multi-indices of a fixed degree grows polynomially.
Each term is a continuous
homogeneous polynomial in $u$, proving the assertion.  This argument
only concerns fixed local target coordinates; no linear structure on
$X$ is being presumed.

\paragraph{The nonlinear holomorphic-map equation.}
The fiber coordinates now have the necessary Sobolev regularity.
It remains to correct their antiholomorphic derivative in the source
while keeping the center and the prescribed first-order directions.
For $u\in W^{1,p_0}(\D,\C^n)$ close to zero, the Sobolev chain rule
in a target chart gives
\[
 \bar\partial\bigl(E(z,u(z))\bigr)
       =E_\xi(z,u(z))\bar\partial u
          +\bar\partial_zE(z,u(z)).
\]
There is no $\bar\partial\overline u$ term, since $E$ is holomorphic
in $\xi$.  Thus the map is weakly holomorphic precisely when
\begin{equation}\label{var:nonlinear-equation}
 \bar\partial u+Q(\cdot,u)=0,
 \qquad
 Q(z,\xi)=E_\xi(z,\xi)^{-1}\bar\partial_zE(z,\xi).
\end{equation}
The formula is independent of the chosen target chart, because the
derivative of a holomorphic chart change multiplies both terms.

Differentiating $\Phi(z,E(z,\xi))=\xi$ with $\xi$ fixed and using
$E_\xi=(D_y\Phi)^{-1}$ gives
\begin{equation}\label{var:q-forward}
 Q(z,\xi)=-\bar\partial_z\Phi(z,E(z,\xi))
       =-\sum_i(\bar\partial\rho_i)(z)C_i(z,E(z,\xi)).
\end{equation}
For the second equality we used holomorphy in $z$ of $C_i$ with $y$
fixed.  In particular $Q$ is holomorphic in $\xi$ and jointly
continuous and bounded on the closed disc times a smaller closed
fiber polydisc.  Cauchy's formula gives coefficients satisfying
$\|Q_\alpha\|_\infty\le Cr_2^{-|\alpha|}$ for some $r_2>0$.
The normally convergent series $\sum_\alpha Q_\alpha u^\alpha$
therefore shows that
\begin{equation}\label{var:q-analytic}
 \mathcal Q\colon W^{1,p_0}(\D,\C^n)\longrightarrow L^{p_0}(\D,\C^n),
 \qquad \mathcal Q(u)=Q(\cdot,u(\cdot)),
\end{equation}
is holomorphic near zero, using the embedding of $W^{1,p_0}$ into
the sup norm.

The constant and linear terms vanish.  Indeed,
$C_i(z,f(z))=0$ and
$D_yC_i(z,f(z))e(z)=I$.  Equations~\eqref{var:inverse-normalization}
and~\eqref{var:q-forward}, together with
$\sum_i\bar\partial\rho_i=0$, therefore give
\begin{equation}\label{var:q-vanishing}
 Q(z,0)=0,\qquad Q_\xi(z,0)=0.
\end{equation}
Equivalently, $\mathcal Q(0)=0$ and $D\mathcal Q(0)=0$ as maps
between the indicated Banach spaces.

\paragraph{The implicit function theorem with the center fixed.}
On the full disc, let $T\colon L^{p_0}\to W^{1,p_0}$ be the solid Cauchy
right inverse from Equation~\eqref{var:cauchy-estimates}.  Since
point evaluation is bounded on $W^{1,p_0}$, the operator
\[
 Hg=Tg-(Tg)(0)
\]
is bounded and satisfies
\begin{equation}\label{var:centered-inverse}
 \bar\partial Hg=g,\qquad (Hg)(0)=0.
\end{equation}
Seek a solution of Equation~\eqref{var:nonlinear-equation} in the
form
\begin{equation}\label{var:u-ansatz}
 u(t,z)=zS(z)t+Hg(t)(z),\qquad g(t)\in L^{p_0}(\D,\C^n).
\end{equation}
Since $S$ is holomorphic and $W^{1,p_0}$, the equation becomes
\[
 G(g,t):=g+\mathcal Q\bigl(zSt+Hg\bigr)=0.
\]
This is a holomorphic equation with values in $L^{p_0}$, and
Equation~\eqref{var:q-vanishing} gives
\[
 G(0,0)=0,\qquad D_gG(0,0)=I.
\]
The complex Banach implicit function
theorem~\cite[Theorem~2.3(ii), (d) of the author's preprint]{Glockner2006}
supplies a holomorphic
solution $g(t)$ for $t$ in a neighborhood of zero.  Differentiating
the equation at zero gives
\[
 g(0)=0,\qquad D_tg(0)=0.
\]
Consequently $u(t)$ is holomorphic with values in $W^{1,p_0}$, is small
in the sup norm for $t$ sufficiently small, and satisfies
\[
 u(t,0)=0,\qquad u(0,z)=0,\qquad
 \frac{\partial u}{\partial t_j}(0,z)=zS(z)e_j,
\]
where $e_j$ here denotes the $j$th standard vector of $\C^n$.

Set $F(t,z)=E(z,u(t,z))$.  The identities in
Equation~\eqref{var:spray-identities} follow from
Equation~\eqref{var:inverse-normalization} and $eS=B$.
Equation~\eqref{var:inverse-substitution} gives its holomorphic
parameter dependence in local $W^{1,p_0}$ coordinates.  In particular
it is jointly continuous up to the source boundary.  For each fixed
$t$, Equation~\eqref{var:nonlinear-equation} makes its coordinate
functions weakly holomorphic on the open disc.  They are therefore
ordinary holomorphic functions.  For each fixed $z$, bounded point
evaluation on $W^{1,p_0}$ and holomorphy of $E$ in its fiber variable
give holomorphy in $t$.  Local separate holomorphy now gives joint
holomorphy on $\{|t|<\eta\}\times\D$.

Finally $F(t,\cdot)$ is $W^{1,p_0}$ with $p_0>2$ and has compact image,
so its area with respect to any smooth Hermitian metric on $X$ is
finite.  Thus all parameters in a sufficiently small complex
neighborhood give admissible centered discs.  The construction imposes
no condition on their boundary maps and requires no extension across
$\partial\D$.
\end{proof}

\begin{remark}\label{var:center-derivative}
Although derivative evaluation at a point is not bounded on the full
space $W^{1,p_0}$, it is bounded on its holomorphic subspace on each
smaller disc, by the Cauchy formula.  In particular
$t\mapsto \partial_zF(t,0)$ is holomorphic.  Hence the central
derivative functional can be differentiated along
Proposition~\ref{var:spray} without additional boundary regularity.
\end{remark}

\section{The trace of the complex Hessian}\label{sec:hessian}

The variations constructed in Section~\ref{var:section} let us apply
the second-derivative test at a maximizing disc. We compute the
central, logarithmically weighted, and unweighted terms separately,
then combine them in one inequality used by both positivity arguments.

Let $f$ be a maximizing disc furnished by Proposition~\ref{ex:maximizer}.
Write $x=f(0)$ and retain the K\"ahler metric $h$ used in its functional.
Use Proposition~\ref{var:frame} to choose a boundary-unitary holomorphic
frame $B=(b_1,\ldots,b_n)$ for $f^*T^{1,0}X$, and
Proposition~\ref{var:spray} to realize the sections $zb_j$ by a family
$F(t,z)$ with fixed center. The arguments in this section apply more
generally to any disc and family with those properties. After decreasing
the Sobolev exponent if necessary, all the maps and frames have
$W^{1,p_0}$ regularity for one $p_0>2$.

Set
\[
 N(z)=B(z)^*h(f(z))B(z).
\]
This positive-definite matrix is smooth in the open disc, continuous
and uniformly positive definite on its closure, and belongs to
$W^{1,p_0}$. Its boundary value is $I$.

\begin{lemma}[Determinant identity]\label{hes:determinant}
With these conventions,
\begin{equation}\label{eq:determinant-identity}
 \log\det N(0)=P_{R_h}(f).
\end{equation}
\end{lemma}

\begin{proof}
In a holomorphic target chart, the matrix of $B$ is holomorphic and
nonsingular. Thus
\[
 \partial_z\partial_{\bar z}\log\det N
 =\partial_z\partial_{\bar z}\log\det h(f(z))
 =-R_h(f',f').
\]
The first equality follows because $\log|\det B|^2$ is harmonic; the
formula is independent of the target chart. The right-hand side lies
in $L^{p_0/2}$. Apply Lemma~\ref{pre:green} and use
$\log\det N=0$ on the boundary.
\end{proof}

For a functional of the family, define
\[
 \cL=\left.\sum_{j=1}^n
   \frac{\partial^2}{\partial t_j\partial\overline{t_j}}\right|_{t=0}.
\]

\begin{proposition}[Trace identities]\label{hes:trace}
For every smooth closed real $(1,1)$-form $q$,
\begin{align}
 \cL d(F_t)&=\tr(B(0)^*k(x)B(0)),\label{eq:hessian-central}\\
 \cL P_q(F_t)&=\frac1{2\pi}\int_0^{2\pi}
      (\tr_h q)(f(e^{i\theta}))\,\dd\theta,\label{eq:hessian-poisson}\\
 \cL A_h(F_t)&=n-A_{R_h}(f).\label{eq:hessian-area}
\end{align}
Here $F_t$ means the disc $z\mapsto F(t,z)$, not a parameter derivative.
\end{proposition}

\begin{proof}
\textbf{The central term and differentiation under the integrals.}
Since $F(t,0)=x$, the vector $\partial_z F(t,0)$ is holomorphic in $t$
and lies in the fixed vector space $T_x^{1,0}X$. Its $t_j$ derivative
at zero is $b_j(0)$. Differentiating its squared $k(x)$ norm gives
Equation~\eqref{eq:hessian-central}. The holomorphic dependence of this
central derivative also follows directly from Cauchy's formula on a
smaller source circle.

We next justify the variations under the area and Poisson integrals.
Choose finitely many fixed source patches and target charts containing
the image of $F$ for sufficiently small $t$. In each chart the map
$t\mapsto F(t,\cdot)$ is holomorphic into $W^{1,p_0}$. The smooth metric
coefficients and the product rule imply that
\[
 t\longmapsto q(F_z(t,\cdot),F_z(t,\cdot))
\]
is twice continuously differentiable into $L^{p_0/2}$; the same holds
with $h$. Indeed, each parameter derivative is a sum of products of two
first source derivatives in $L^{p_0}$, with uniformly bounded Sobolev
coefficient factors. A fixed partition of unity in the source gives the
global assertion. Since $p_0/2>1$, integration against $1$ or against
$\log(1/|z|)$ is a continuous functional on this space. Differentiation
under both integrals is therefore valid.

\paragraph{The local identity and its Poisson integral.}
Locally write $q=\ddc\phi$. Joint holomorphy of $F$ on the open product
of the parameter domain and the disc gives
\begin{align*}
 q(F_z,F_z)&=\partial_z\partial_{\bar z}(\phi\circ F),\\
 \cL(\phi\circ F)&=\sum_j q(F_{t_j},F_{t_j})\big|_{t=0}.
\end{align*}
Commuting these derivatives and using $F_{t_j}(0,z)=zb_j(z)$ yields the
intrinsic identity
\begin{equation}\label{eq:local-hessian}
 \cL q(F_z,F_z)
 =\partial_z\partial_{\bar z}
     \left(|z|^2\tr(B^*q(f)B)\right).
\end{equation}
In particular, the mixed source derivative on the right belongs to
$L^{p_0/2}$ by the preceding differentiability argument.

The scalar function in parentheses is continuous on $\overline\D$ and
vanishes at zero. Lemma~\ref{pre:green} therefore computes its Poisson
integral from its boundary values. On the boundary, the columns of $B$
are an $h$-orthonormal basis, so
$\tr(B^*qB)=\tr_h q$. This proves
Equation~\eqref{eq:hessian-poisson}.

\paragraph{The area integral and its weak boundary limit.}
At the identity matrix, trace and log determinant have the same
differential. The boundary normalization will therefore let us replace
one radial derivative by the other, with an error controlled by the
radial Sobolev estimate along suitable radii.
For the area identity put $T=\tr N$. Integrating
Equation~\eqref{eq:local-hessian} with $q=h$ on $|z|<r$ gives
\begin{equation}\label{eq:area-radius-hessian}
 \frac1\pi\int_{|z|<r}\cL h(F_z,F_z)\,\dd S
 =r^2\langle T\rangle_r+\frac{r^3}{2}\langle T_r\rangle_r.
\end{equation}
The left-hand side tends to $\cL A_h(F_t)$ by integrability. The first
term on the right tends to $n$. To handle the second term without
assuming boundary differentiability, observe that
\begin{equation}\label{eq:trace-logdet-error}
 T_r-(\log\det N)_r=\tr((I-N^{-1})N_r).
\end{equation}
All matrix norms here are equivalent, and uniform positive definiteness
gives
\[
 \big|\langle T_r-(\log\det N)_r\rangle_r\big|
 \le C\|N(r,\cdot)-I\|_{L^2(S^1)}
       \|N_r(r,\cdot)\|_{L^2(S^1)}.
\]
The radial $W^{1,2}$ estimate and the trace $N=I$ give
\begin{equation}\label{eq:radial-trace-estimate}
 \|N(r,\cdot)-I\|_2
 \le(1-r)^{1/2}
       \left(\int_r^1\|N_s(s,\cdot)\|_2^2\,\dd s\right)^{1/2}
 =o((1-r)^{1/2}).
\end{equation}
This follows first on almost every radial line by the fundamental
theorem for Sobolev functions, and then by Cauchy--Schwarz and
integration over the angle. Smoothness inside the disc and continuity
at the boundary give the same estimate for the interior radii used
below.

Since $\int_{1/2}^1\|N_r\|_2^2\,\dd r<\infty$, there are radii
$r_\nu\uparrow1$ for which
$(1-r_\nu)\|N_r(r_\nu,\cdot)\|_2^2$ is bounded. Otherwise its eventual
lower bound would contradict integrability. Along these radii,
Equations~\eqref{eq:trace-logdet-error} and
\eqref{eq:radial-trace-estimate} show that the two circular derivative
means differ by $o(1)$.

On the other hand, the curvature calculation in
Lemma~\ref{hes:determinant} and Equation~\eqref{eq:green-area} give, for
every interior radius,
\[
 \frac r2\langle(\log\det N)_r\rangle_r=-A_{R_h,r}(f).
\]
The right-hand side tends to $-A_{R_h}(f)$. Substitute into
Equation~\eqref{eq:area-radius-hessian} along the chosen radii. Its
left-hand side has a limit independent of the sequence, proving
Equation~\eqref{eq:hessian-area}.
\end{proof}

\begin{corollary}[Maximizing inequality]\label{hes:maximum}
For a maximizer of $d+P_q-\delta A_h$ centered at $x$, the
boundary-unitary frame satisfies
\begin{equation}\label{eq:maximum-inequality}
 \tr(B(0)^*k(x)B(0))
 +\frac1{2\pi}\int_0^{2\pi}(\tr_hq)(f(e^{i\theta}))\,\dd\theta
 \le\delta\bigl(n-A_{R_h}(f)\bigr).
\end{equation}
It also satisfies Equation~\eqref{eq:determinant-identity}, and
$A_q(f)\ge-d(f)$.
\end{corollary}

\begin{proof}
Every sufficiently small parameter value gives an admissible
finite-area disc with the same center. Hence the real-valued twice
differentiable function $J(F_t)$ has a local maximum at zero, and
$\cL J(F_t)\le0$. Proposition~\ref{hes:trace} gives the displayed
inequality. The remaining assertions are Lemma~\ref{hes:determinant}
and the dilation inequality of Lemma~\ref{ex:dilation}.
\end{proof}

The two applications use different parts of the maximizing property.
In Section~\ref{sec:nefness}, we set $h=k$ and $q=R_k$ and let
$\delta$ tend to zero after comparing the maximum with each fixed test
disc; this gives a bound on $P_{R_k}$ for all finite-area discs.
In Section~\ref{sec:volume}, we set $q=-h$ and $\delta=1$.
Comparison with the constant disc then bounds $P_h(f)+A_h(f)$ by
$D$, independently of $h$. Together with $R_h\ge-h$, this is the
additional information that makes the volume estimate uniform over
all such metrics.

\section{Bounded canonical potentials and nefness}
\label{sec:nefness}

The first application of the maximizing inequality proves nefness of
$K_X$. We proceed in three steps: bound the Poisson Ricci integral on
every finite-area disc,
convert the bound into locally bounded plurisubharmonic metric weights,
and smooth those weights with arbitrarily small loss of positivity.
We first apply the preceding disc inequalities with the fixed K\"ahler
form $k$.  As before, $R_k=-\ddc\log\det k$, and $D$ denotes the
uniform upper bound for $d(f)=k(f'(0),f'(0))$ from
Lemma~\ref{ex:brody}.

\begin{proposition}[Uniform Ricci-disc bound]
\label{nef:disc-bound}
There is a constant $C\geq 0$, depending only on the fixed metric $k$,
such that
\begin{equation}
 P_{R_k}(f)\leq C
 \label{eq:nef-uniform-disc-bound}
\end{equation}
for every finite-area holomorphic disc $f:\D\to X$.
\end{proposition}

\begin{proof}
For $0<\delta\leq 1$ and a prescribed center $x$, let $f$ maximize
\[
J_\delta(f)=d(f)+P_{R_k}(f)-\delta A_k(f).
\]
Such a maximizer exists by Proposition~\ref{ex:maximizer}.
Use the boundary-normalized holomorphic frame $B$ associated to this
maximizer and put $N(z)=B(z)^*k(f(z))B(z)$.  Set
\[
 M=\sup_X|\tr_k R_k|.
\]
The dilation inequality, Equation~\eqref{eq:tail}, gives
$A_{R_k}(f)\geq-d(f)\geq-D$.
Consequently, Equation~\eqref{eq:maximum-inequality}, with $h=k$ and
$q=R_k$, implies
\[
 \tr N(0)
 \leq \delta\bigl(n-A_{R_k}(f)\bigr)+M
 \leq n+D+M.
\]
The determinant identity, Equation~\eqref{eq:determinant-identity},
and the arithmetic-geometric mean inequality yield
\[
 P_{R_k}(f)=\log\det N(0)
 \leq n\log\!\left(\frac{n+D+M}{n}\right)=:C_1.
\]
Since $A_k(f)\geq0$ and $d(f)\leq D$, the maximal value of
$J_\delta$ is at most $D+C_1$, independently of both $x$ and $\delta$.
Now fix any finite-area disc $g$, use its center in this maximization,
and obtain
\[
 d(g)+P_{R_k}(g)-\delta A_k(g)\leq D+C_1.
\]
Letting $\delta\downarrow0$ for this fixed disc and using $d(g)\geq0$
proves the assertion with $C=D+C_1$.  This limiting argument requires
no area bound uniform in $\delta$ for the maximizing discs.
\end{proof}

\subsection{The Poisson envelope on the determinant bundle}

Poletsky introduced disc envelopes for plurisubharmonic functions,
and Rosay established the manifold case \cite{Poletsky1991,Rosay2003}.
The formulation below is due to Drinovec Drnov\v sek and Forstneri\v c,
whose theorem also covers locally irreducible complex spaces
\cite{DrinovecForstneric2012}. The use of a line bundle lets the
Ricci-disc bound control these envelopes without choosing one global
potential on $X$.

We use the Poletsky--Rosay Theorem in the following form.  If $Z$ is
a connected complex manifold and $b:Z\to\R\cup\{-\infty\}$ is upper
semicontinuous, then
\[
 \mathcal P b(y)
 =\inf_{G(0)=y}\frac{1}{2\pi}\int_0^{2\pi}b(G(e^{i\theta}))\,\dd\theta
\]
is plurisubharmonic or identically $-\infty$.  Here the infimum is
taken over maps $G:\overline{\D}\to Z$ that are continuous on
$\overline{\D}$ and holomorphic in $\D$.  There is no requirement that
$Z$ be compact or that $b$ be bounded below.  This is, in particular,
the manifold case of the more general Theorem~1.1 of
\cite{DrinovecForstneric2012}.  For continuous real-valued $b$, the
same envelope results if the test maps are required to be holomorphic
on neighborhoods of $\overline{\D}$: replace each test map $G$ by
$z\mapsto G(rz)$ and let $r\uparrow1$.

\begin{proposition}[Bounded positive canonical potentials]
\label{nef:bounded-potentials}
There is a bounded real-valued function $u$ on $X$ such that
\[
 T=-R_k+\ddc u\geq0
\]
is a closed positive $(1,1)$-current.  In every local holomorphic
frame $s_i$ of $\det T^{1,0}X$, put $a_i=\log|s_i|_k^2$.
Then $a_i+u$ is plurisubharmonic and locally bounded.  These are
local weights of a singular Hermitian metric with nonnegative
curvature on $K_X$.
\end{proposition}

\begin{proof}
Write $L=\det T^{1,0}X$ and let $\pi:Y=L\setminus 0_X\to X$ be
the complement of its zero section.  Define
\[
 \ell(s)=\log|s|_k^2,\qquad s\in Y.
\]
This is a smooth finite-valued function on the connected complex
manifold $Y$.  In a local trivialization $s=\lambda s_i(x)$ it has
the expression
\[
 \ell(s)=\log|\lambda|^2+a_i(x),\qquad
 \ddc a_i=-R_k.
\]
Since $\lambda\ne0$, the first summand is pluriharmonic.  Therefore
\begin{equation}
 \ddc\ell=-\pi^*R_k.
 \label{eq:nef-total-space-curvature}
\end{equation}

Let $v=\mathcal P\ell$, using test discs holomorphic past the closed
unit disc.  For such a disc $G$ with $G(0)=s$, its projection
$f=\pi\circ G$ has finite area.  Green's formula and
Equation~\eqref{eq:nef-total-space-curvature} give
\[
 \frac{1}{2\pi}\int_0^{2\pi}\ell(G(e^{i\theta}))\,\dd\theta
 =\ell(s)-P_{R_k}(f).
\]
Proposition~\ref{nef:disc-bound} bounds every such mean below by
$\ell(s)-C$, while the constant disc has mean $\ell(s)$.  Thus
\begin{equation}
 \ell-C\leq v\leq\ell.
 \label{eq:nef-envelope-bounds}
\end{equation}
The Poletsky--Rosay Theorem applies even though $\ell$ is globally
unbounded, and Equation~\eqref{eq:nef-envelope-bounds} excludes its
identically $-\infty$ alternative.  Hence $v$ is plurisubharmonic.

For every $\lambda\in\C^*$, multiplication by $\lambda$ is a
biholomorphism of $Y$ and gives a bijection between the test discs
centered at $s$ and at $\lambda s$.  Since
$\ell(\lambda s)=\ell(s)+\log|\lambda|^2$, it follows that
\[
 v(\lambda s)=v(s)+\log|\lambda|^2.
\]
Thus $v-\ell$ is constant on every fiber.  It has the form
$u\circ\pi$, where $-C\leq u\leq0$.  Pulling back $v$ by the local
holomorphic map $s_i$ proves that $a_i+u$ is plurisubharmonic.
It is locally bounded because $a_i$ is smooth and $u$ is bounded.
The local currents $\ddc(a_i+u)$ agree and equal $-R_k+\ddc u$.

To check the line-bundle sign explicitly, if $s_i=g_{ij}s_j$ then
\[
 a_i-a_j=\log|g_{ij}|^2.
\]
Let $s_i^*$ be the dual frame of $K_X=L^*$.  The prescription
\[
 |s_i^*|^2=\exp\bigl(-(a_i+u)\bigr)
\]
is compatible with $s_i^*=g_{ij}^{-1}s_j^*$.  Its curvature is
$\ddc(a_i+u)=T\geq0$.  In particular $T$ represents
$2\pi c_1(K_X)$, rather than $2\pi c_1(L)$.
\end{proof}

\subsection{From bounded potentials to nefness}

A holomorphic line bundle is nef if, for every $\eps>0$, it admits a
smooth Hermitian metric whose real curvature is at least $-\eps k$.
This definition is independent of the background K\"ahler metric on
a compact manifold. The following approximation result converts the
bounded potentials just constructed into such metrics.

\begin{lemma}[Smoothing a bounded positive potential]
\label{nef:bounded-smoothing}
Let $(M,k)$ be a compact K\"ahler manifold, let $\alpha$ be a smooth
closed real $(1,1)$-form on $M$, and let $u$ be a bounded real-valued
function such that $\alpha+\ddc u\geq0$ in the sense of currents.
For every $\eps>0$ there is a smooth real-valued function $u_\eps$
on $M$ such that
\[
 \alpha+\ddc u_\eps\geq-\eps k.
\]
\end{lemma}

This is the bounded-potential case of the regularization theorem for
positive currents: locally bounded plurisubharmonic potentials have
zero Lelong numbers, so \cite[Corollary~6.4]{Demailly1992} gives
nefness. Appendix~\ref{app:smoothing} supplies a direct proof. Its
key estimate compares regularizations on overlapping charts without
assuming that the original bounded potentials are continuous.

\begin{proposition}[Nefness of the canonical class]
\label{nef:nef}
For every $\eps>0$ there is a smooth real-valued function $u_\eps$
on $X$ such that
\[
 -R_k+\ddc u_\eps\geq-\eps k.
\]
Consequently $2\pi c_1(K_X)$ is nef.
\end{proposition}

\begin{proof}
Apply Lemma~\ref{nef:bounded-smoothing} with $M=X$, $\alpha=-R_k$,
and the bounded function $u$ of Proposition~\ref{nef:bounded-potentials}.
The resulting functions $u_\eps$ give smooth canonical metrics with
local squared norms $\exp(-(a_i+u_\eps))$ in the dual frames $s_i^*$.
Their curvatures are $-R_k+\ddc u_\eps\geq-\eps k$.
\end{proof}

\section{A uniform volume estimate}
\label{sec:volume}

Keep the background K\"ahler form $k$ fixed, and let $D$ be the
uniform bound from Lemma~\ref{ex:brody}:
\[
 d(f)=k\bigl(f'(0),f'(0)\bigr)\le D
\]
for every holomorphic disc in $X$.  We first extract a second consequence
of the extremal-disc argument.  Its constant will be independent of the
K\"ahler metric to which it is applied.

\begin{proposition}[Uniform lower bound for the volume form]
\label{prop:volume-lower-bound}
For every smooth K\"ahler form $h$ on $X$ satisfying
$R_h\ge-h$, one has
\begin{equation}
\label{eq:uniform-volume-lower-bound}
 h^n\ge c_{n,D}\,k^n,
 \qquad
 c_{n,D}=e^{-D}\left(\frac{n}{2n+D}\right)^n>0.
\end{equation}
\end{proposition}

\begin{proof}
Fix $x\in X$.  Use the maximizer from Proposition~\ref{ex:maximizer} with
$q=-h$ and $\delta=1$, so that its functional is
\[
 J(f)=d(f)-P_h(f)-A_h(f).
\]
The constant disc centered at $x$ has value zero.  Consequently a
maximizing disc $f$ satisfies
\begin{equation}
\label{eq:volume-extremal-energy}
 P_h(f)+A_h(f)\le d(f)\le D.
\end{equation}
Both terms on the left are nonnegative.  The assumed Ricci inequality
therefore gives
\[
 A_{R_h}(f)\ge-A_h(f)\ge-D,
 \qquad
 P_{R_h}(f)\ge-P_h(f)\ge-D.
\]

Let $B$ be the holomorphic frame along $f$ normalized to be
$h$-orthonormal on the boundary.  Since $\tr_h(-h)=-n$, the
maximum inequality, Equation~\eqref{eq:maximum-inequality}, yields
\[
 \tr\bigl(B(0)^*k(x)B(0)\bigr)
 \le 2n-A_{R_h}(f)\le2n+D.
\]
Meanwhile, the determinant identity,
Equation~\eqref{eq:determinant-identity}, gives
\[
 \det\bigl(B(0)^*h(x)B(0)\bigr)
 =\exp\bigl(P_{R_h}(f)\bigr)\ge e^{-D}.
\]
The arithmetic--geometric mean inequality for the eigenvalues of the
positive Hermitian matrix $B(0)^*k(x)B(0)$ gives
\[
 \det\bigl(B(0)^*k(x)B(0)\bigr)
 \le\left(\frac{2n+D}{n}\right)^n.
\]
Taking the quotient cancels the common factor $|\det B(0)|^2$:
\[
 \frac{h^n}{k^n}(x)
 =\frac{\det\bigl(B(0)^*h(x)B(0)\bigr)}
        {\det\bigl(B(0)^*k(x)B(0)\bigr)}
 \ge e^{-D}\left(\frac{n}{2n+D}\right)^n.
\]
The point $x$ was arbitrary, proving the assertion.
\end{proof}

The maximizing disc, its boundary regularity, and the size of its
holomorphic variation may all depend on $h$ and $x$.  The proof applies
the variational identities separately for each such choice.  None of
these auxiliary bounds enters the constant in
Equation~\eqref{eq:uniform-volume-lower-bound}.

\section{Positive top self-intersection}
\label{sec:positive-intersection}

Nefness allows us to approach the canonical class through K\"ahler
classes. We choose in each such class a metric to which the uniform
volume estimate applies. This Monge--Amp\`ere and cohomological-volume
strategy is also used by Wu and Yau
\cite[Proposition~8(i) and the proof of Lemma~6 in the authors' preprint]{WuYau2016}.
We use the following established form of the negative-sign complex
Monge--Amp\`ere theorem.  It is the Aubin--Yau existence theorem
\cite{Aubin1998,Yau1978}; the statement for an arbitrary smooth volume form is
also recorded explicitly in \cite[\S1, Equation~(1.1)]{Berman2019}.

\begin{theorem}[Negative-sign complex Monge--Amp\`ere theorem]
\label{thm:negative-cma}
Let $M$ be a compact K\"ahler manifold of dimension $n$, let $\gamma$
be a K\"ahler form on $M$, and let $\Omega$ be a smooth positive volume
form.  There is a smooth real function $w$ such that
\[
 \gamma+\ddc w>0,
 \qquad
 (\gamma+\ddc w)^n=e^w\Omega.
\]
\end{theorem}

There is no prescribed normalization of $\Omega$ in this statement.
The unknown includes its additive constant, which enters the
right-hand side through $e^w$.  In particular, the theorem is not
restricted to a K\"ahler form representing the canonical class.

\begin{proposition}
\label{prop:positive-top-intersection}
The canonical class has positive top self-intersection:
\[
 \int_X c_1(K_X)^n>0.
\]
\end{proposition}

\begin{proof}
Write
\[
 \alpha=[-R_k]=2\pi c_1(K_X).
\]
The nefness proved in Proposition~\ref{nef:nef} means that, for each
$t>0$, there is
a smooth real function $a_t$ for which
\[
 -R_k+\ddc a_t\ge-\frac t2 k.
\]
It follows that
\[
 \gamma_t=-R_k+tk+\ddc a_t\ge\frac t2 k>0
\]
is a K\"ahler form.  Apply Theorem~\ref{thm:negative-cma} with this
form and with the smooth positive volume form
$\Omega_t=e^{a_t}k^n$.  The resulting K\"ahler form
\[
 h_t=\gamma_t+\ddc w_t
\]
satisfies
\begin{equation}
\label{eq:twisted-cma-volume}
 h_t^n=e^{a_t+w_t}k^n.
\end{equation}
Taking $-\ddc$ of the logarithm of the ratio of the volume forms in
Equation~\eqref{eq:twisted-cma-volume} shows, with our Ricci convention, that
\begin{equation}
\label{eq:twisted-ricci-equation}
 R_{h_t}=R_k-\ddc(a_t+w_t)=-h_t+tk\ge-h_t.
\end{equation}
Proposition~\ref{prop:volume-lower-bound} therefore applies, with the
same constant for every $t>0$:
\[
 \int_Xh_t^n\ge c_{n,D}\int_Xk^n.
\]
On the other hand $[h_t]=\alpha+t[k]$, so closedness and Stokes'
theorem give
\[
 \int_Xh_t^n=\int_X(\alpha+t[k])^n.
\]
The expression on the right is a polynomial in $t$.  Passing to its
value at zero yields
\begin{equation}
\label{eq:positive-canonical-volume}
 (2\pi)^n\int_Xc_1(K_X)^n
 =\int_X\alpha^n
 \ge c_{n,D}\int_Xk^n>0.
\end{equation}
This passage uses only cohomology.  No convergence of the metrics
$h_t$, or uniform bounds on their potentials, is required.
\end{proof}

\section{Bigness, projectivity, and ampleness}
\label{sec:ampleness}

At this point $K_X$ is nef and has positive top self-intersection.
We first obtain enough sections to make $X$ Moishezon, then use its
K\"ahler structure to obtain projectivity. Only after that step do
we invoke the projective cone theorem to prove ampleness.
We give the positivity argument in the form needed here, including the
estimate that permits the use of holomorphic Morse inequalities.
For a real curvature form $\Theta$ of a smooth Hermitian line-bundle
metric, let $X(\Theta,\le1)$ denote the locus where $\Theta$ is
nondegenerate and has at most one negative eigenvalue.  The number of
negative eigenvalues is independent of the chosen background
Hermitian metric.

Demailly's strong holomorphic Morse inequalities
\cite[Theorem~0.1 and Equation~(0.7)]{Demailly1985} imply the following
criterion: on a compact complex manifold, if
\begin{equation}
\label{eq:morse-bigness-criterion}
 \int_{X(\Theta,\le1)}\Theta^n>0,
\end{equation}
then the line bundle is big and the manifold is Moishezon
\cite[Theorem~0.8(a)]{Demailly1985}.  Here ``big'' means that the
line bundle has Kodaira dimension $n$; equivalently its spaces of
sections have growth of order $m^n$ along sufficiently divisible
tensor powers.  A compact complex manifold is Moishezon when its field
of meromorphic functions has transcendence degree $n$.

\begin{lemma}[Nefness and positive intersection imply bigness]
\label{lem:nef-big}
Let $L$ be a nef holomorphic line bundle on a compact connected K\"ahler
manifold $(X,k)$ of dimension $n\ge1$.  If
$\int_Xc_1(L)^n>0$, then $L$ is big and $X$ is Moishezon.
\end{lemma}

\begin{proof}
Put $\beta=2\pi c_1(L)$ and $V=\int_X\beta^n>0$.
For $0<\eps\le1$, choose a smooth metric on $L$ with real curvature
$\Theta_\eps\ge-\eps k$.  Thus
\[
 P_\eps=\Theta_\eps+2\eps k
\]
is positive definite.  At a point where $\Theta_\eps$ has a negative
eigenvalue, write its eigenvalues relative to $k$ as
$\lambda_1,\ldots,\lambda_n$ and choose $j$ with $\lambda_j<0$.
Then
\[
 |\lambda_j|\le\eps,
 \qquad
 |\lambda_i|\le\lambda_i+2\eps\quad(i\ne j).
\]
Consequently
\[
 \left|\prod_{i=1}^n\lambda_i\right|
 \le\eps\prod_{i\ne j}(\lambda_i+2\eps)
 \le\eps\sum_{j=1}^n\prod_{i\ne j}(\lambda_i+2\eps).
\]
In terms of volume forms this is
\begin{equation}
\label{eq:negative-index-mass}
 |\Theta_\eps^n|\le n\eps\,P_\eps^{n-1}\wedge k
\end{equation}
on the entire locus with a negative eigenvalue.  The right-hand side
has integral
\[
 n\eps\int_X(\beta+2\eps[k])^{n-1}\wedge[k]=O(\eps).
\]
The implied constant is independent of the chosen metric: the
remaining cohomological integral is a fixed polynomial in $\eps$,
bounded for $0\le\eps\le1$.

The degenerate locus contributes zero to $\Theta_\eps^n$.  Removing
the loci of index at least two from its total integral therefore
changes that integral by at most $O(\eps)$ in absolute value.  Since
$\int_X\Theta_\eps^n=V$, we obtain
\[
 \int_{X(\Theta_\eps,\le1)}\Theta_\eps^n
 \ge V-O(\eps)>0
\]
for a sufficiently small fixed $\eps$.  The criterion in
Equation~\eqref{eq:morse-bigness-criterion} applies to this one smooth metric
and proves the assertion.  In dimension one the locus of index at
least two is empty; the same argument, with $P_\eps^0=1$, applies
without modification.
\end{proof}

Apply Lemma~\ref{lem:nef-big} to $L=K_X$, using Proposition~\ref{nef:nef} and
Proposition~\ref{prop:positive-top-intersection}.  We conclude that
$K_X$ is big and $X$ is Moishezon.  The following established
projectivity theorem now applies.

\begin{theorem}[Moishezon's projectivity theorem]
\label{thm:moishezon-projectivity}
A compact Moishezon manifold that admits a K\"ahler metric is
projective.
\end{theorem}

For a primary source proving a stronger form of this theorem, see
Theorem~6 in the author's preprint of \cite{Namikawa2002}: it permits
Moishezon spaces with
$1$-rational singularities.  A smooth manifold satisfies that
hypothesis.  Thus $X$ is now a smooth projective variety, and we may
apply the projective cone theorem.

The following standard implication is recorded in
\cite[Lemma~2.1]{DiverioTrapani2019}. We give the argument through an
effective perturbation of the canonical divisor and the log cone
theorem.

\begin{proposition}[Bigness in the absence of rational curves]
\label{prop:no-rational-ample}
Let $X$ be a smooth projective variety over $\C$.  If $K_X$ is big
and $X$ contains no rational curves, then $K_X$ is ample.
\end{proposition}

\begin{proof}
Kodaira's lemma gives an integer $m>0$, an ample Cartier divisor $H$,
and an effective divisor $E$ such that
\begin{equation}
\label{eq:kodaira-decomposition}
 mK_X\sim H+E.
\end{equation}
For the needed form of this lemma, choose a smooth very ample
hyperplane section $H$. Bigness gives $h^0(X,mK_X)\ge cm^n$ along
sufficiently divisible powers, whereas the weak holomorphic Morse
inequality gives $h^0(H,mK_X|_H)=O(m^{n-1})$
\cite[Theorem~0.1(a)]{Demailly1985}. The restriction sequence therefore
gives a nonzero section of $mK_X-H$ for some $m$, producing $E$.
For $n=1$, $H$ is a finite set and the same estimate is $O(1)$.

Choose a sufficiently small positive rational number $\eta$ so that
$(X,\eta E)$ is Kawamata log terminal.  This choice is possible even
when $E$ is singular or nonreduced. Choose a log resolution
$\pi:Y\to X$ of $(X,E)$, whose existence follows from
\cite[Theorem~35]{Kollar2007}, and write
\[
 K_Y-\pi^*K_X=\sum_i a_iF_i,
 \qquad \pi^*E=\widetilde E+\sum_i b_iF_i,
\]
where $F_i$ are exceptional prime divisors and $\widetilde E$ is the
strict transform with its multiplicities. The discrepancy criterion
for a Kawamata log terminal pair is given in
\cite[\S4.4]{Fujino2011}. The exceptional discrepancies have the form
$a_i-\eta b_i$, with $a_i\ge0$ because $X$ is smooth and $b_i\ge0$
because $E$ is effective.  There are finitely many such coefficients.
Taking $\eta>0$ sufficiently small keeps each discrepancy greater
than $-1$ and each strict-transform boundary coefficient less than
one, which is the required condition.

The log cone theorem says that a projective Kawamata log terminal
pair $(X,\Delta)$ over $\C$ with effective rational boundary has a
rational curve of negative $(K_X+\Delta)$-degree whenever
$K_X+\Delta$ is not nef.  This follows, for example, from the cone
decomposition and rational-curve assertion in
\cite[Theorem~1.1, especially part~(5)]{Fujino2011}; for a Kawamata
log terminal pair the non-log-canonical locus in that statement is
empty.  The theorem has no restriction on the dimension.

Since $X$ has no rational curves, it follows that
\[
 N=K_X+\eta E
\]
is nef, in the numerical sense that its degree on every curve is
nonnegative. Combining this with Equation~\eqref{eq:kodaira-decomposition} gives
\[
 (1+\eta m)K_X\sim_{\mathbb Q}N+\eta H.
\]
The sum of a nef rational divisor and an ample rational divisor on a
projective variety is ample by Kleiman's ampleness criterion
\cite[Definition~4.9 and Theorem~4.10]{Fujino2011}: it is positive
on every nonzero element of the closed cone of curves. Hence
$(1+\eta m)K_X$ is ample. Clearing the
positive rational coefficient shows that a positive integral multiple
of $K_X$ is ample, and therefore $K_X$ itself is ample.
\end{proof}

\begin{proof}[Completion of the proof of Theorem~\ref{thm:main}]
A rational curve in $X$ would give a nonconstant holomorphic map
$\mathbb P^1\to X$ through its normalization.  Its restriction to
$\C\subset\mathbb P^1$ would still be nonconstant, by the identity
theorem.  The hypothesis therefore excludes rational curves.
We have proved that $X$ is projective and $K_X$ is big, so
Proposition~\ref{prop:no-rational-ample} proves that $K_X$ is ample.
Equivalently, a sufficiently high positive tensor power of $K_X$ is
very ample and its global holomorphic sections define a holomorphic
embedding into projective space.  This is the assertion in every
positive complex dimension.
\end{proof}

\section{Two consequences}\label{sec:consequences}

We now combine canonical ampleness with two separate results. The first
gives an explicit exponent for global generation of the pluricanonical
bundles. The second describes a finite cover when the universal cover
has an additional semialgebraic realization. Neither companion result
is used in the proof of Theorem~\ref{thm:main}.

\begin{corollary}[Pluricanonical freeness]\label{cor:pluricanonical-freeness}
Let $X$ be a compact connected Brody-hyperbolic K\"ahler manifold of
complex dimension $n\ge1$. Then $K_X^{\otimes m}$ is generated by its
global sections for every integer $m\ge n+2$.
\end{corollary}
\begin{proof}
Theorem~\ref{thm:main} makes $X$ smooth projective with $K_X$ ample.
Apply Fujita freeness in its all-powers form
\cite[Corollary~6.3]{OpenAIFujita2026} with $L=K_X$ and exponent
$m-1\ge n+1$; the resulting globally generated adjoint bundle is
$K_X\otimes L^{\otimes(m-1)}=K_X^{\otimes m}$.
\end{proof}

\begin{corollary}[Finite-cover product]\label{cor:semialgebraic-product}
Let $X$ be a compact connected Brody-hyperbolic K\"ahler manifold of
positive complex dimension. Suppose its ordinary universal cover is
biholomorphic to a semialgebraic open subset of a complex projective
variety, with openness in the ordinary complex topology. Then $X$
admits a connected finite \'etale cover biholomorphic to
\[
                    (D/\Gamma_0)\times F,
\]
where $D$ is a bounded symmetric domain, $\Gamma_0$ is a discrete group
of biholomorphisms acting freely, properly discontinuously and
cocompactly on $D$, and $F$ is a simply connected compact
Brody-hyperbolic projective manifold. Either $D$ or $F$ may be a point.
\end{corollary}
\begin{proof}
Theorem~\ref{thm:main} makes $X$ projective. The semialgebraic-cover
classification \cite[Theorem~1.1]{OpenAISemialgebraicCovers2026} gives
$\widetilde X\simeq D\times\C^m\times F$, with $F$ simply connected,
normal and projective. Smoothness of this product forces $F$ to be
smooth. Brody hyperbolicity passes to covers and factors, so $m=0$
and $F$ is Brody hyperbolic. If $\dim F>0$, Theorem~\ref{thm:main}
makes $K_F$ ample, and Kobayashi's finite-automorphism theorem
\cite[Theorem, p.~184]{Kobayashi1959} shows that
$\operatorname{Aut}(F)$ is finite. This also holds when $F$ is a point.

For a deck transformation of $D\times F$, the first coordinate is
constant on each compact connected fiber $F$. Applying the same
observation to its inverse shows that its
base map is an automorphism of $D$ and its fiber maps are
automorphisms of $F$. Since $D$ is connected and
$\operatorname{Aut}(F)$ is finite, these fiber maps are independent
of the base point. Thus the deck group $\Gamma$ acts by products.
The kernel $\Gamma_0$ of its homomorphism to
$\operatorname{Aut}(F)$ has finite index and acts faithfully and
freely on $D$. Proper discontinuity on $D$ follows by applying the
deck-action property to $K\times F$ for compact $K\subset D$; in
particular $\Gamma_0$ is discrete. Hence
$\widetilde X/\Gamma_0\simeq(D/\Gamma_0)\times F$ is a connected
finite unramified holomorphic cover of $X$, and therefore a finite
\'etale cover of the projective manifold $X$. Its compactness implies
that $D/\Gamma_0$ is compact.
\end{proof}



\appendix
\section{Smoothing a bounded positive potential}
\label{app:smoothing}

We prove Lemma~\ref{nef:bounded-smoothing}. The problem is to smooth
a bounded potential while losing arbitrarily little positivity.
Bounded plurisubharmonic functions need not be continuous, so uniform
convergence of their regularizations is unavailable. Instead, we
compare regularizations at fixed multiples of the same small radius.
The comparison becomes uniformly small and permits gluing by a
regularized maximum.

\begin{lemma}[Fixed-factor changes of radius]
\label{nef:radius-comparison}
Let $\psi$ be plurisubharmonic on a domain in $\C^n$.  Suppose that
all closed balls $\overline B(x,R)$ with centers in a fixed set $E$
are contained in the domain, and that $m\leq\psi\leq M$ on their
union.  Put
\[
 S_x(r)=\sup_{|z-x|\leq r}\psi(z),\qquad 0<r\leq R.
\]
For every $A>1$ and $Ar<R$,
\begin{equation}
 0\leq S_x(Ar)-S_x(r)
 \leq\frac{(M-m)\log A}{\log(R/r)},\qquad x\in E.
 \label{eq:nef-radius-comparison}
\end{equation}
For each sufficiently small fixed $r$, the function $x\mapsto S_x(r)$
is plurisubharmonic on its open domain of definition.
\end{lemma}

\begin{proof}
First, $t\mapsto S_x(e^t)$ is nondecreasing and convex.  To see
convexity, fix $0<r_1<r_2\leq R$ and compare the restriction of
$\psi$ to each complex line through $x$ with the harmonic function
on its annulus $r_1<|z-x|<r_2$ that is affine in $\log|z-x|$ and
takes the constant values $S_x(r_1)$ and $S_x(r_2)$ on the boundary
circles.  The subharmonic maximum principle bounds $\psi$ by this
function on each annulus.  On the inner ball the bound follows from
the monotonicity of the interpolation.  Taking the supremum over an
intermediate ball proves convexity.

Apply this convexity at the three logarithmic radii $\log r$,
$\log r+\log A$, and $\log R$.  It gives
\[
 S_x(Ar)-S_x(r)
 \leq\frac{\log A}{\log(R/r)}\bigl(S_x(R)-S_x(r)\bigr),
\]
which implies Equation~\eqref{eq:nef-radius-comparison}.

For fixed $r$, the supremum over the closed translation ball is
upper semicontinuous.  Indeed, for a convergent sequence of centers,
choose maximizing translation vectors and pass to a convergent
subsequence in $\overline B(0,r)$; upper semicontinuity of $\psi$
gives the required upper bound.  It is also a locally bounded
supremum of the plurisubharmonic translations
$x\mapsto\psi(x+\zeta)$, $|\zeta|\leq r$.  Its upper
semicontinuity therefore makes it plurisubharmonic.
\end{proof}

\begin{proof}[Proof of Lemma~\ref{nef:bounded-smoothing}]
\textbf{Local potentials and regularization.}
Choose finitely many coordinate patches $U_i$, smooth real functions
$a_i$ with $\alpha=\ddc a_i$ on neighborhoods of $\overline U_i$,
and open sets
\[
 V_i\Subset W_i\Subset U_i
\]
such that the $V_i$ cover $M$. Such local potentials exist because
$\alpha$ is a smooth closed real $(1,1)$-form. On overlaps,
$a_i-a_j$ is pluriharmonic.

The distributional positivity of $\alpha+\ddc u$ means that
$a_i+u$ has a unique plurisubharmonic representative $\psi_i$.
These representatives are locally bounded: their almost-everywhere
bounds follow from boundedness of $u$ and smoothness of $a_i$, and
the submean inequality and upper semicontinuity extend these bounds
to every point. Moreover,
\[
 \psi_i-\psi_j=a_i-a_j
\]
pointwise on overlaps, since the two plurisubharmonic functions
$\psi_i$ and $\psi_j+a_i-a_j$ agree almost everywhere.
Thus, after replacing $u$ on a set of measure zero, we may write
$\psi_i=a_i+u$ everywhere. The functions $\psi_i$ are bounded above
and below on $U_i$, with bounds uniform over the finite collection.
The positive distances between the compact sets $\overline W_i$
and the boundaries of $U_i$ allow all subsequent regularizations
to be defined on neighborhoods of $\overline W_i$.

In the coordinates of $U_i$, let
\[
 S_{i,b}(x)=\sup_{|y-x|\leq b}\psi_i(y).
\]
Fix a smooth nonnegative kernel of integral one supported in the
unit ball, and let $\rho_{b/4}$ be its rescaling to radius $b/4$.
For all sufficiently small $b>0$, set
\[
 \varphi_{i,b}=S_{i,b}*\rho_{b/4}
\]
near $\overline W_i$.  Lemma~\ref{nef:radius-comparison} and
convolution show that $\varphi_{i,b}$ is smooth and
plurisubharmonic.  If $|y-x|\leq b/4$, then
\[
 B(x,b/2)\subset B(y,b)\subset B(x,2b).
\]
Taking suprema and then averaging proves
\begin{equation}
 S_{i,b/2}(x)\leq\varphi_{i,b}(x)\leq S_{i,2b}(x).
 \label{eq:nef-convolution-bounds}
\end{equation}

\paragraph{Comparison on overlaps.}
We next prove the uniform overlap comparison
\begin{equation}
 \sup_{W_i\cap W_j}
 \left|(\varphi_{i,b}-a_i)-(\varphi_{j,b}-a_j)\right|
 \longrightarrow0\quad\text{as }b\downarrow0.
 \label{eq:nef-overlap-comparison}
\end{equation}
All coordinate changes and their inverses have bounded derivatives
on neighborhoods of the compact overlap sets
$\overline W_i\cap\overline W_j$.  After decreasing the allowed
radius, coordinate balls centered on these overlaps are therefore
comparable by fixed factors.  Thus there are constants $c,C>0$,
independent of the center and of $b$, such that an $i$-coordinate
ball of radius $b$ contains the $j$-coordinate ball of radius $cb$
and is contained in the $j$-coordinate ball of radius $Cb$.
The balls in question lie in $U_i\cap U_j$.

The identity $\psi_i-\psi_j=a_i-a_j$ holds on this intersection,
and the smooth function $a_i-a_j$ varies by $O(b)$ across any such
ball, uniformly in the center.  These facts, together with
Equation~\eqref{eq:nef-convolution-bounds}, give constants
$c_1,c_2,C_2>0$ for which
\[
 S_{j,c_1b}(x)-a_j(x)-C_2b
 \leq\varphi_{i,b}(x)-a_i(x)
 \leq S_{j,c_2b}(x)-a_j(x)+C_2b.
\]
The analogous bounds for $\varphi_{j,b}-a_j$ have radii $b/2$
and $2b$.  Lemma~\ref{nef:radius-comparison}, using a fixed radius
available near each compact overlap, compares all four suprema.
It bounds their differences by $O(1/|\log b|)$.  This proves
Equation~\eqref{eq:nef-overlap-comparison}, in fact with an error
bounded by $O(1/|\log b|)+O(b)$.  All constants can be chosen
uniformly because there are only finitely many pairs of patches.

\paragraph{Gluing the local functions.}
Choose smooth functions $\chi_i$ on neighborhoods of $\overline W_i$
such that
\[
 -3\leq\chi_i\leq0,\qquad
 \chi_i=0\text{ on }V_i,\qquad
 \chi_i=-3\text{ near }\partial W_i.
\]
There is a fixed constant $C_0\geq0$ satisfying
$\ddc\chi_i\geq-C_0k$ for every $i$.  Given $\sigma>0$, choose
$b$ so small that every discrepancy in
Equation~\eqref{eq:nef-overlap-comparison} is less than $\sigma$,
and set on $W_i$
\[
 f_i=\varphi_{i,b}-a_i+\sigma\chi_i.
\]

Here is a precise regularized maximum with the compatibility needed
when the set of available patches changes.  Fix a smooth even
probability density $\eta_\sigma$ on $\R$, supported in
$(-\sigma/4,\sigma/4)$.  For each nonempty finite index set $I$
define
\[
 \mathcal M_I((t_i)_{i\in I})
 =\int_{\R^I}\max_{i\in I}(t_i+s_i)
       \prod_{i\in I}\eta_\sigma(s_i)\,\dd s_i.
\]
This is smooth, convex, nondecreasing in every argument, and
satisfies
\begin{equation}
 \mathcal M_I(t+c\mathbf1)=\mathcal M_I(t)+c.
 \label{eq:nef-max-translation}
\end{equation}
If $t_i$ is more than $\sigma/2$ below the maximum of the other
arguments, then it never attains the shifted maximum in the
integrand.  Its removal leaves exactly $\mathcal M_{I\setminus\{i\}}$,
because the unused probability density integrates to one.

For $x\in M$, put $I(x)=\{i:x\in W_i\}$ and define
\[
 u_\sigma(x)=\mathcal M_{I(x)}((f_i(x))_{i\in I(x)}).
\]
This is globally smooth.  To verify the only issue, fix
$x_0\in\partial W_i$ and choose $j$ with $x_0\in V_j$.
On a neighborhood of $x_0$ within $W_i$, we have
$\chi_i=-3$ and $\chi_j=0$, so the overlap comparison implies
\[
 f_i<f_j-2\sigma.
\]
Thus branch $i$ is inactive throughout this neighborhood, including
under all the allowed shifts.  The exact deletion property just
proved removes it from the formula before crossing $\partial W_i$.
There are finitely many indices, so this argument removes all
indices whose patch boundaries pass through $x_0$, leaving a smooth
formula on a neighborhood of $x_0$.

\paragraph{The curvature estimate.}
Work near any point in a fixed
coordinate patch with local potential $a_j$ for $\alpha$.  Each active branch has
\[
 a_j+f_i=\varphi_{i,b}+(a_j-a_i)+\sigma\chi_i.
\]
The difference $a_j-a_i$ is pluriharmonic, since both functions
are local potentials for $\alpha$. Hence
\begin{equation}
 \ddc(a_j+f_i)\geq-C_0\sigma k.
 \label{eq:nef-branch-curvature}
\end{equation}
By Equation~\eqref{eq:nef-max-translation},
$a_j+u_\sigma$ is the regularized maximum of these branches.
The first derivatives of $\mathcal M_I$ are nonnegative and sum
to one; its Hessian is positive semidefinite.  The complex Hessian
chain rule therefore expresses $\ddc(a_j+u_\sigma)$ as a convex
combination of the branch Hessians plus a nonnegative $(1,1)$-form.
Equation~\eqref{eq:nef-branch-curvature} consequently gives
\[
 \alpha+\ddc u_\sigma=\ddc(a_j+u_\sigma)
 \geq-C_0\sigma k.
\]
Choose $\sigma>0$ with $C_0\sigma\leq\eps$, and rename this
function $u_\eps$. The resulting smooth closed form represents
$[\alpha]$ and has the required lower bound.
\end{proof}


\begingroup
\raggedright
\small
\bibliographystyle{plain}
\bibliography{references}
\endgroup
\end{document}
